% Generated by maint/texbuild/theory_tex.py from theory/workbooks/cell_tracking/cell_tracking_table.md.
% Edits here are overwritten. Edit the markdown, then run the script again.
\chapter{\texorpdfstring{The cell tracking table}{The cell tracking table}}
\label{chap:cell-tracking-table}

\textbf{Purpose:} The micro n-body problem cell\_\allowbreak{}tracking exists to solve, part by part. For each part: the physics we hold it to, what the program does about it today, what it wants to do, every hypothesis tried there with its result, and the next move.

\textbf{Scope:} the program: \texttt{cell\_\allowbreak{}tracking/\allowbreak{}src/\allowbreak{}}, the rules in \texttt{cell\_\allowbreak{}tracking/\allowbreak{}base.\allowbreak{}cfg}, the cell-specific modules still under \texttt{engine/\allowbreak{}} until they move (\texttt{bodies}, \texttt{score\_\allowbreak{}sample}, \texttt{answer\_\allowbreak{}key}, \texttt{measure/\allowbreak{}*}, \texttt{group\_\allowbreak{}objects}, \texttt{link\_\allowbreak{}objects}, \texttt{relate\_\allowbreak{}frames}), the physics in the thought experiments (the harmonics source above all, \texttt{noise\_\allowbreak{}sieve\_\allowbreak{}5\_\allowbreak{}cell\_\allowbreak{}tracking\_\allowbreak{}harmonics.\allowbreak{}pdf}, in \texttt{thought\_\allowbreak{}experiments/\allowbreak{}cell\_\allowbreak{}tracking/\allowbreak{}}; the rest are in \texttt{thought\_\allowbreak{}experiments/\allowbreak{}engine/\allowbreak{}}), and the numbers in \hyperref[chap:ledger]{ledger.md}. The machine the program runs on is a separate concern with its own table, engine\_\allowbreak{}table.md; a row here names the machine rows (M1 to M12) it reads. Statuses follow \hyperref[chap:readme]{README.md}. A hypothesis with no run in the ledger says so.

\section{\texorpdfstring{Two tables, two concerns}{Two tables, two concerns}}

The engine is the machine: exact arithmetic, the residual operator, the component tree, the moments, the one correlation operator, the marginal, the record machine, the files. The program is cell tracking: what a cell is, which body it is next frame, when it divides or dies. Each table is optimized on its own terms. They are optimized against each other only through the interface at the end of each table (what this program asks of the engine, what the engine offers). Merged, neither can be optimized on its own (Doug, 23 September). A second program, RSNA knee MRI, now runs on the same engine, which is why the machine must stay free of cells.

\section{\texorpdfstring{The problem}{The problem}}

A developing zebrafish tissue is an n-body system: cells are the bodies, suspended in a fluid medium, touching and adhering through their membranes, moving, dividing, dying, entering and leaving the volume. We see it only as 100 frames of a 64×256×256 volume, with each sample's own deterministic noise on top. From that we must recover every body, its state, its motion, which body it is in the next frame, and its lineage. Everything must be exact, and no number may be chosen.

The score pays for three things: nodes matched to key nodes within 7 µm, edges between matched nodes, and divisions (0.1 of the score). Every row below is judged by what it moves on that score.

\section{\texorpdfstring{The driver: scan everything, sort at the end}{The driver: scan everything, sort at the end}}

Doug, 25 September: “you scan everything and then you sort stuff. You don't need to sort it before you scan in all of the samples.” Today's driver partitions every frame as it reads it: the cut at ℓ* (row 1) fixes each voxel's body from one frame before the sample's other 99 frames are seen. That is the moment deciding membership, which row 1's and row 9's physics forbid. The score reads only node positions (within 7 µm) and edges, never regions. So the new driver first scans every frame of every sample into points and fields that bound nothing. Only then does it sort: link, keep, divide, and grade.

A peak here is a point: a positive voxel of the residual that no one of its 26 neighbours exceeds. Nothing is assigned to it. The regions around peaks (the watershed's catchment regions, found by tobogganing, each voxel following its steepest neighbour) bound the set, and the driver does not form them. A body's shape is read from the field at its peak, not from a region: the residual's exact second differences there are the ℓ = 2 of its local field, with its long axis and aspect. Everything in this table is tracking and lives in \texttt{cell\_\allowbreak{}tracking/\allowbreak{}}. The engine is called only for math: the residual, the drift, the entropy history, and one-to-one matching. Every run and every test is a job on the tessera schedule; nothing touches the device outside it.

Doug, 25 September: “The whole point of the tracker is to track cells. If you can't walk back something you did you built it wrong.” Every stage keeps what it came from, and every decision keeps what it replaced:

\begin{itemize}
\item a point keeps its frame, voxel and exact level;
\item a transform keeps its exact inverse;
\item a link keeps its two points and its evidence;
\item a choice in the sort keeps the options it passed over, with their weights.
\end{itemize}

Nothing is deleted. What the sort does not keep is marked, and every mark can be walked back to the readings. A map that sends two readings to one (a clip, a rounding, a bin) cannot be walked back and is not built.

The organism physics each stage answers to is the thought experiments' (the harmonics source, \texttt{noise\_\allowbreak{}sieve\_\allowbreak{}5\_\allowbreak{}cell\_\allowbreak{}tracking\_\allowbreak{}harmonics.\allowbreak{}pdf}, and the two fluidic ones), row by row in the table below:

\begin{itemize}
\item the gradients define themselves, and nothing is tuned;
\item a body's identity is the limit of its coherence over the whole sample, not a tag;
\item every body's ℓ ≤ 2 harmonics are its own, from intrinsic differences (mass density, membrane elasticity, aspect ratio), and oblate bodies make ℓ = 1 and ℓ = 2 dominate;
\item the field moves as a whole (drift) and each body moves on top of it; a body's velocity is its dipole's time derivative, its next place is predicted from that before anything is matched, and v\_\allowbreak{}origin = −∇(∂C₁ₘ/∂t) points back to where the motion started;
\item bodies touch and adhere through their membranes, move together, and two membranes cannot pass through one another;
\item a split is a boundary discontinuity; a division's daughters part along the parent's long axis with its mass conserved, and the local entropy bumps and settles;
\item a lysis is a bump that rises into the floor and does not settle;
\item mitosis and lysis change the fluidics for a great distance around them: a disturbance runs out through the tissue, overlapping disturbances are told apart by their harmonics, and each is traced back to its origin;
\item oblate bodies deforming, snapping or throwing a pseudopod shed predictable dipole waves along their axis;
\item what decays to maximum entropy is the floor; what holds low and moves is a body;
\item every event is judged over the whole sample, never in the moment.
\end{itemize}

\tablerecord{\textbf{S0. The set's intensities}}
\begin{description}
\item[{what it does}] Every sample is read once, and the set's readings are counted into one exact histogram. Each sample gets its own affine footing, (reading − lo, hi − lo), with lo and hi converged from their 1\% and 99\%. The contrast is the set's cumulative count. Both are walked back exactly, and nothing is clipped (O2). Each sample's counts and footing go into its \texttt{.\allowbreak{}readings}, and the set's total, width and samples into \texttt{scan.\allowbreak{}set}. The count itself, C, goes into every \texttt{.\allowbreak{}points}, so each level walks back from its own file
\item[{organism physics (row)}] contrast across the set; the medium stays E1's to subtract (3)
\item[{over}] the whole set, before the scan
\item[{reads}] \texttt{engine\_\allowbreak{}iapx\_\allowbreak{}load}
\item[{lives in}] \texttt{cell\_\allowbreak{}tracking/\allowbreak{}src/\allowbreak{}scan/\allowbreak{}} (\texttt{scan\_\allowbreak{}set}, its first pass)
\item[{status}] built 25 September. \textbf{Measured} on the 25: 10,485,760,000 readings, so the contrast is 34 bits. Every footing converged in 6 to 16 steps, lo 34 to 699 and hi 330 to 2,099. Every \texttt{.\allowbreak{}points} holds the C rebuilt from the 25 \texttt{.\allowbreak{}readings}, checked by a separate reader on 25 of 25
\end{description}

\tablerecord{\textbf{S1. Peaks}}
\begin{description}
\item[{what it does}] every positive local maximum of the residual, as a point (t, z, y, x) with its exact level and the faces that cut its 26-neighbourhood (O1, O4)
\item[{organism physics (row)}] a body stands above the medium (1); the medium is subtracted whole (3)
\item[{over}] each frame, as it is read
\item[{reads}] \texttt{engine\_\allowbreak{}residual\_\allowbreak{}planes} (E1): S0's contrast as planes of its width, with the residual's width returned
\item[{lives in}] \texttt{cell\_\allowbreak{}tracking/\allowbreak{}src/\allowbreak{}peaks/\allowbreak{}}, run by \texttt{cell\_\allowbreak{}tracking/\allowbreak{}src/\allowbreak{}scan/\allowbreak{}}
\item[{status}] built 25 September; peaks at 1, 9, 10 and 12 limbs \textbf{proved} against an exact host scan (131 checks, 0 failed). \textbf{Measured} on the 25: 6,518,615 points, 1,145 to 6,031 a frame, and 942,418 of them (14.5\%) have a face in their neighbourhood. Each level is 10 limbs. The residual takes 20.1 ms a frame and the peaks 0.86 ms. A separate reader of every \texttt{.\allowbreak{}points} finds 0 levels not positive, 0 voxels out of order and 0 outside the frame. The faces are counted, not yet stored per point (O4)
\end{description}

\tablerecord{\textbf{S1b. Shape}}
\begin{description}
\item[{what it does}] at each peak, the residual's six exact second differences: the local field's ℓ = 2, its long axis and aspect
\item[{organism physics (row)}] the body's own harmonics (2); oblate bodies (2, 9)
\item[{over}] each frame, as it is read
\item[{reads}] the frame's residual
\item[{lives in}] \texttt{cell\_\allowbreak{}tracking/\allowbreak{}src/\allowbreak{}peaks/\allowbreak{}}
\item[{status}] not built
\end{description}

\tablerecord{\textbf{S2. Drift}}
\begin{description}
\item[{what it does}] the lag that best carries frame t's positive set onto t+1's
\item[{organism physics (row)}] the field moves as a whole first (4)
\item[{over}] each frame pair, as it is read
\item[{reads}] the engine's drift (E3, A4)
\item[{lives in}] the engine; called from the scan
\item[{status}] built in the engine; the scan does not call it yet
\end{description}

\tablerecord{\textbf{S3. History}}
\begin{description}
\item[{what it does}] per voxel and bit, the flips over the whole sample
\item[{organism physics (row)}] the floor and what holds low (12); the bumps of division and lysis (9, 10)
\item[{over}] the whole sample
\item[{reads}] \texttt{engine\_\allowbreak{}entropy\_\allowbreak{}set} (M8)
\item[{lives in}] the engine; \texttt{-\allowbreak{}-\allowbreak{}run entropy}
\item[{status}] built
\end{description}

\tablerecord{\textbf{S4. The set and the null}}
\begin{description}
\item[{what it does}] each sample's peaks in its null frames (row 8's draws), and what the set's samples say together, read before any sample is sorted
\item[{organism physics (row)}] a body is real where it stands above the field-noise reading (8)
\item[{over}] every sample
\item[{reads}] S1 to S3
\item[{lives in}] \texttt{cell\_\allowbreak{}tracking/\allowbreak{}src/\allowbreak{}}
\item[{status}] theory
\end{description}

\tablerecord{\textbf{S5. Motion, then links}}
\begin{description}
\item[{what it does}] each point's next place is predicted first: the drift, plus its own velocity (its last link's displacement, the dipole's time derivative), plus the local flow (its neighbours' departures from the drift, weighted by level). Its candidates are the gate (O6): its nearest point of t+1, and every point of t+1 whose nearest prediction it is, with ties kept. The cost is the weighted squared distance from the prediction, ‖·‖²\_\allowbreak{}w with w = σ∘σ, weighed by the sample's own count of first-rank links at that distance (O7), plus the distance between the two shapes (S1b, O8), the shape taken from the chain's last point so a fate change is not a break (H2). The first pass is one to one by \texttt{heaviest\_\allowbreak{}matching} on exact integers. Two neighbours' links may not cross: the sign of (c\_\allowbreak{}A − c\_\allowbreak{}B)·(c\_\allowbreak{}A′ − c\_\allowbreak{}B′) is kept
\item[{organism physics (row)}] motion (5); local flow (4); identity by harmonics (2, 7); membranes do not cross (6)
\item[{over}] the whole sample, after its scan
\item[{reads}] S1, S1b, S2; the engine's matching (M7)
\item[{lives in}] \texttt{cell\_\allowbreak{}tracking/\allowbreak{}src/\allowbreak{}}
\item[{status}] not built
\end{description}

\tablerecord{\textbf{S6. Chains and identity}}
\begin{description}
\item[{what it does}] Linked points form chains, solved once over the whole sample as a min-cost flow (O10). Every cost is an exact integer on one footing: links from O7; “no link” from the null draw; appearance free at a face or the first frame. The local moves (O12), repeated to a fixed point, check the solve. A chain is kept where it stands above the same chains in the null frames, over the whole sample, with no length chosen (O17). A chain that ends and one that starts at its predicted place, any number of frames on, are one chain: the gap is bridged (O14). Chains that part and rejoin, or a point ending beside one starting, are one body with two tops (O5, O13, O15). Each link's marginal is taken over its support component (O18), and each source keeps its heaviest outcome (O19)
\item[{organism physics (row)}] identity is coherence (7); the null (8)
\item[{over}] the whole sample
\item[{reads}] S5, S4
\item[{lives in}] \texttt{cell\_\allowbreak{}tracking/\allowbreak{}src/\allowbreak{}}
\item[{status}] theory
\end{description}

\tablerecord{\textbf{S7. Division}}
\begin{description}
\item[{what it does}] A chain goes on as two: two points in t+1 whose nearest prediction is the parent's (O6), parting along the parent's long axis (S1b), with the history bumping at t and settling (S3). This is the evidence O9 learns. The daughters need not be alike: mass is conserved on their sum, and each takes its own harmonics (H1). The split sits at the frame of the strongest evidence along the chain, and a daughter chain that does not stand above the null is no daughter (O16). The neighbours' departures from the drift in the frames after point back to it
\item[{organism physics (row)}] a split is a discontinuity; mass and axis; the bump that settles; the disturbance (9, 4)
\item[{over}] the whole sample
\item[{reads}] S5, S6, S1b, S3
\item[{lives in}] \texttt{cell\_\allowbreak{}tracking/\allowbreak{}src/\allowbreak{}}
\item[{status}] theory
\end{description}

\tablerecord{\textbf{S8. Lysis}}
\begin{description}
\item[{what it does}] a chain ends where the history bumps and does not settle, and its neighbours' departures point back to it; a chain that ends with no bump is a missed detection, not a death
\item[{organism physics (row)}] the bump that stays; the disturbance (10, 4)
\item[{over}] the whole sample
\item[{reads}] S6, S3
\item[{lives in}] \texttt{cell\_\allowbreak{}tracking/\allowbreak{}src/\allowbreak{}}
\item[{status}] theory
\end{description}

\tablerecord{\textbf{S9. Faces}}
\begin{description}
\item[{what it does}] a chain that starts or ends where its prediction leaves the view entered or left, exactly, with no buffer (O4, O11)
\item[{organism physics (row)}] entering and leaving (11)
\item[{over}] the whole sample
\item[{reads}] S5, S6
\item[{lives in}] \texttt{cell\_\allowbreak{}tracking/\allowbreak{}src/\allowbreak{}}
\item[{status}] theory
\end{description}

\tablerecord{\textbf{S10. Disturbances}}
\begin{description}
\item[{what it does}] every division and lysis (S7, S8) is an origin; the departures around it over the following frames are resolved by their ℓ = 1 and ℓ = 2 about each origin, so overlapping disturbances are told apart and traced back; they then correct S5's local flow, and the links are sorted again
\item[{organism physics (row)}] mitosis and lysis change the fluidics far around them; dipole waves; retrograde tracing (4)
\item[{over}] the whole sample
\item[{reads}] S5, S7, S8
\item[{lives in}] \texttt{cell\_\allowbreak{}tracking/\allowbreak{}src/\allowbreak{}}
\item[{status}] theory
\end{description}

\tablerecord{\textbf{S11. Output}}
\begin{description}
\item[{what it does}] the kept chains' points as nodes, consecutive points as edges, and every division's parent link; the graph's invariants checked first, with the unexplained ends and starts counted per frame (O20, H4); written for \texttt{maint/\allowbreak{}score\_\allowbreak{}submission.\allowbreak{}py} and as the submission CSV
\item[{organism physics (row)}] the score (7 µm, edges, divisions)
\item[{over}] each sample
\item[{reads}] S6 to S9
\item[{lives in}] \texttt{cell\_\allowbreak{}tracking/\allowbreak{}src/\allowbreak{}}
\item[{status}] not built
\end{description}

Order of building: S0 (O2's norm, once the set is ingested); S1 (with O4's faces), S1b and S2 in the scan; then S5 with O6's gate, and S11 with every point kept (the first scored run); then O7's weights, S4, and S6 as O10's one solve with O12 to O15; then S7, S8 and S9 (O9, O16, O11); then S10 with O21's flow. O18 and O20 come in beside S6 and S11.

\subsection{\texorpdfstring{OrganoidTracker's techniques, into the driver}{OrganoidTracker's techniques, into the driver}}

These were read on 25 September from \texttt{vendor\_\allowbreak{}for\_\allowbreak{}ref/\allowbreak{}Organoid\allowbreak{}Tracker/\allowbreak{}} (read only), following \texttt{organoid\_\allowbreak{}tracker\_\allowbreak{}predict\_\allowbreak{}positions.\allowbreak{}py}, \texttt{\_\allowbreak{}predict\_\allowbreak{}divisions.\allowbreak{}py}, \texttt{\_\allowbreak{}create\_\allowbreak{}tracks.\allowbreak{}py}, \texttt{\_\allowbreak{}marginalization.\allowbreak{}py} and \texttt{\_\allowbreak{}callibrate\_\allowbreak{}marginalization.\allowbreak{}py} into the modules they call. Their pipeline runs in this order:

\begin{enumerate}
\item A network calls positions.
\item Nearest-neighbour candidates are drawn.
\item Networks score each link, and each position's chance of dividing.
\item One min-cost flow runs over the whole movie (dpct).
\item Local repairs fix what the flow missed.
\item Each link gets a marginal, and errors are flagged.
\end{enumerate}

Each technique is listed with its chosen numbers, then our form of it: exact, nothing chosen, and in the sort after the scan.

\tablerecord{\textbf{O1. Peak calling}}
\begin{description}
\item[{OrganoidTracker (where)}] \texttt{position\_\allowbreak{}predictor.\allowbreak{}py} (\texttt{peak\_\allowbreak{}local\_\allowbreak{}max}); \texttt{loss\_\allowbreak{}functions.\allowbreak{}peak\_\allowbreak{}finding}
\item[{what it does there}] A voxel is a peak where it equals the maximum over a (3, 13, 13) box and exceeds 0.1 of the range. At inference: min\_\allowbreak{}distance 6 px and threshold 0.1, after interpolating 5 layers between z planes (\texttt{reconstruct\_\allowbreak{}volume}), so peaks are called on a near-isotropic grid. Every voxel that ties its box maximum is a peak.
\item[{ours: exact, nothing chosen}] S1 as built: the 26 neighbours, no box, and no threshold past > 0. The z anisotropy enters the sort's norm, w = σ∘σ, not interpolated planes. A tie goes to the lower index.
\item[{stage}] S1
\item[{status}] built
\end{description}

\tablerecord{\textbf{O2. Intensity normalization}}
\begin{description}
\item[{OrganoidTracker (where)}] \texttt{position\_\allowbreak{}predictor.\allowbreak{}\_\allowbreak{}split\_\allowbreak{}into\_\allowbreak{}patches} (the division and link predictors do the same)
\item[{what it does there}] Per frame, the readings are mapped onto [0, 1] between the 1\% and 99\% quantiles taken over the network's time window, then clipped. An optional top-hat (a Gaussian blur, then a minimum and a maximum filter of \texttt{tophat\_\allowbreak{}mask}) is subtracted.
\item[{ours: exact, nothing chosen}] Taken (Doug, 25 September), for contrast, once the set is ingested. Every sample is read once, and the set's readings are counted into one exact histogram, one count per 16-bit value. The bounds lo and hi are its quantiles: each is the least value whose cumulative count reaches its share of the set, exactly. Their clip is not built, because it cannot be walked back: every reading past a bound becomes the bound. What is built comes in two parts, both walked back exactly. (a) \textbf{The footing:} each sample's own affine map, (reading − lo) over (hi − lo), is held as that pair and never divided. E1 is linear and cancels a constant, so this map moves no peak. What it does is put samples of different brightness on one footing for the set's readings (S4). One pair for the whole set would change nothing, since every sample already shares the raw scale, so the pairs are per sample. They are read once the set is ingested. (b) \textbf{The contrast:} the set's cumulative count, v ↦ \#\{readings of the set ≤ v\}, is exact histogram equalization. It rises strictly on every value the set holds, so its inverse is the table of those values. It needs no levels, and it changes the scan's peaks the way the clip meant to. It carries the contrast (Doug, 25 September: “cumulative count sound like more information”). For (a), the levels are not written in (Doug, 25 September: “we can let the numbers converge on them”). They start at theirs, 1\% and 99\%, and iterate to a fixed point on the set. The rule (built 25 September; 6 to 16 steps on the 25) is that each bound is the two-means split of its half of the set: the readings below the median for lo, and those above it for hi. Each is iterated from its seed until it stops moving. The class sums and counts are exact integers, a reading is compared against the midpoint by cross-multiplying, and a one-dimensional two-means stops in finitely many steps. The top-hat is not taken: the residual subtracts the medium (E1).
\item[{stage}] S0
\item[{status}] built
\end{description}

\tablerecord{\textbf{O3. The time window}}
\begin{description}
\item[{OrganoidTracker (where)}] the position, division and link networks
\item[{what it does there}] Each network sees frames t + dt over its window, a stack rather than one frame.
\item[{ours: exact, nothing chosen}] The scan reads one frame at a time and keeps everything. Every judgement that needs other frames is made in the sort (S4 to S10), over the whole sample, with no window.
\item[{stage}] S4 to S10
\item[{status}] the design
\end{description}

\tablerecord{\textbf{O4. The border}}
\begin{description}
\item[{OrganoidTracker (where)}] \texttt{custom\_\allowbreak{}filters.\allowbreak{}get\_\allowbreak{}edges}; \texttt{links\_\allowbreak{}postprocessor.\allowbreak{}postprocess}
\item[{what it does there}] Training ignores voxels within (2.5, 11, 11) px of the x and y faces and the last z plane. Positions within \texttt{margin\_\allowbreak{}xy} of the edge are dropped.
\item[{ours: exact, nothing chosen}] Each point records which faces cut its 26-neighbourhood, exactly and with no margin. S9 reads the record.
\item[{stage}] S1, S9
\item[{status}] not built
\end{description}

\tablerecord{\textbf{O5. A dividing cell called twice}}
\begin{description}
\item[{OrganoidTracker (where)}] \texttt{division\_\allowbreak{}predictor.\allowbreak{}remove\_\allowbreak{}division\_\allowbreak{}oversegmentation}
\item[{what it does there}] Two points both called dividing and within 4.5 µm of each other (1.5 × that if both are dividing) become one point at their midpoint. In the frame before, points within 4.5 µm of a dividing point, other than its nearest, are removed.
\item[{ours: exact, nothing chosen}] Two peaks of one frame stay two points, and the sort decides. Chains that part and go on are a division (S7). Chains that rejoin are one body with two tops (E9's reconvergence, S6). The saddle between two neighbouring peaks is an exact number: the level of their LCA in the component tree (low fruit 6).
\item[{stage}] S6, S7
\item[{status}] theory
\end{description}

\tablerecord{\textbf{O6. The candidate gate}}
\begin{description}
\item[{OrganoidTracker (where)}] \texttt{nearest\_\allowbreak{}neighbor\_\allowbreak{}linker.\allowbreak{}nearest\_\allowbreak{}neighbor}
\item[{what it does there}] Each point of t+1 goes to the points of t within tolerance × its nearest distance, at most 5 of them; each point of t goes to t+1 the same way; the candidates are the union. The link network uses tolerance 2. A point whose place in the other frame is outside the image is skipped.
\item[{ours: exact, nothing chosen}] From each point's prediction (S5): its nearest point of t+1 under ‖·‖²\_\allowbreak{}w. From each point of t+1: its nearest prediction. The gate is the union, with ties kept. Two points of t+1 whose nearest prediction is the same point are that point's division candidates (S7). A prediction outside Ω is a leaving (S9). No tolerance and no count.
\item[{stage}] S5
\item[{status}] not built
\end{description}

\tablerecord{\textbf{O7. The link's weight, learned without a key}}
\begin{description}
\item[{OrganoidTracker (where)}] \texttt{nearest\_\allowbreak{}neighbor\_\allowbreak{}linker.\allowbreak{}nearest\_\allowbreak{}neighbor\_\allowbreak{}probabilities}
\item[{what it does there}] Every point's shortest out-link is taken as true and its others as false. The links are counted in log-distance bins of 0.25, and a logistic curve is fitted to the counts. The penalty is −log₁₀ of the odds.
\item[{ours: exact, nothing chosen}] The same count, with no bins and no fit. ‖·‖²\_\allowbreak{}w from the prediction is an exact integer, so at each value d² the sample's (or the set's) count of first-rank links over all links is an exact rational, P(first | d²). The count needs the whole sample scanned before any link is weighed.
\item[{stage}] S5, S4
\item[{status}] theory
\end{description}

\tablerecord{\textbf{O8. The link network}}
\begin{description}
\item[{OrganoidTracker (where)}] \texttt{link\_\allowbreak{}predictor}
\item[{what it does there}] Its inputs: patches at both ends over the time window, and the displacement given as CoordConv channels. Its output is Platt-scaled against curated links.
\item[{ours: exact, nothing chosen}] What it learns from is the displacement and the two ends' look. S5 reads the displacement from the prediction exactly, and S1b's pair of shapes for the look. Nothing is calibrated.
\item[{stage}] S5, S1b
\item[{status}] not built
\end{description}

\tablerecord{\textbf{O9. The division network}}
\begin{description}
\item[{OrganoidTracker (where)}] \texttt{division\_\allowbreak{}predictor}
\item[{what it does there}] Gives each point a probability from a patch over its window, Platt-scaled; the penalty is −log₁₀ of the odds. It is trained on the mother's last frame (or 3 frames before to 2 after the split), with dying cells shown more often.
\item[{ours: exact, nothing chosen}] The evidence is read from the data: the shape's aspect along its long axis (S1b), the history's bump (S3), and two t+1 points sharing one nearest prediction (O6). The frames around a split are where the bump rises and settles, read from S3, not a window.
\item[{stage}] S7
\item[{status}] theory
\end{description}

\tablerecord{\textbf{O10. The global solve}}
\begin{description}
\item[{OrganoidTracker (where)}] \texttt{dpct\_\allowbreak{}linker} (Haubold, Ales, Wolf, Hamprecht, ECCV 2016: generalized successive shortest paths; or Magnusson)
\item[{what it does there}] One min-cost flow over every frame at once. A node costs 2 if left unused and 0 if used. Appearance and disappearance cost their penalties, and are free at the first and last frames. A node may send two units (a division) where three things hold: its division penalty is under 2, it has more than one future, and its second-best out-link beats its disappearance. Five weights (link, detection, division, appearance, disappearance), all 1. A link is offered only within 4 of the best in and out penalties, and under 4.
\item[{ours: exact, nothing chosen}] S6 is this solve: once over the whole sample, after the scan. Every cost is an exact integer on one footing: a link from O7; a point's “no link” from its null draw (S4: the nearest it finds in a frame where no correspondence exists), as the marginal already weighs “none” (row 7). Appearance is free where the back-prediction leaves Ω (S9) and at the first frame. A division costs by O9's evidence. No weights and no cut-offs: O6's gate is the pruning. The machine's \texttt{heaviest\_\allowbreak{}matching} (M7) solves one frame pair, one to one. The flow over the sample, with two-unit sources, is the tracker's own, in \texttt{cell\_\allowbreak{}tracking/\allowbreak{}src/\allowbreak{}}. Whether a generic min-cost flow belongs in the engine is Doug's call.
\item[{stage}] S6
\item[{status}] theory
\end{description}

\tablerecord{\textbf{O11. Appearance at the faces}}
\begin{description}
\item[{OrganoidTracker (where)}] \texttt{dpct\_\allowbreak{}linker.\allowbreak{}calculate\_\allowbreak{}appearance\_\allowbreak{}penalty}
\item[{what it does there}] Within 8 µm of any face but z = 0, the chance of appearing rises linearly from 0.01 to 0.51. A random ±0.05 is added to every penalty “to help with optimalization”.
\item[{ours: exact, nothing chosen}] S9: a start is an entering exactly when its back-prediction lies outside Ω, with no buffer. Ties break exactly (least index, least ‖·‖²\_\allowbreak{}w), never by noise.
\item[{stage}] S9
\item[{status}] not built
\end{description}

\tablerecord{\textbf{O12. Local moves}}
\begin{description}
\item[{OrganoidTracker (where)}] \texttt{links\_\allowbreak{}postprocessor.\allowbreak{}finetune\_\allowbreak{}solution}, 4 passes
\item[{what it does there}] Drop a link where appearance + disappearance is cheaper. Reconnect a loose start or end by breaking another track, when the total drops. Add a daughter where the division plus the link (plus the other's disappearance) is cheaper. Swap two crossing links. Drop a division whose penalty is over 2 by cutting the worse daughter.
\item[{ours: exact, nothing chosen}] The same moves, as exact integer comparisons of the same costs, repeated until no move lowers the total: a fixed point with no pass count. On an exact whole-sample solve they change nothing, so they are the check that it is optimal, and the repair where the gate cut an option. The swap is the no-crossing rule (S5, row 6) seen from the cost side.
\item[{stage}] S6
\item[{status}] theory
\end{description}

\tablerecord{\textbf{O13. A loop from over-detection}}
\begin{description}
\item[{OrganoidTracker (where)}] \texttt{links\_\allowbreak{}postprocessor.\allowbreak{}connect\_\allowbreak{}loose\_\allowbreak{}ends}
\item[{what it does there}] A track ends and another starts within 4 frames. If the link + 2 < disappearance + appearance, they are joined and the points between are deleted (---===--- → ---------).
\item[{ours: exact, nothing chosen}] A chain that parts and rejoins is one body with two tops over that stretch, whatever its length (E9's reconvergence). It is kept against the null, not a penalty.
\item[{stage}] S6
\item[{status}] theory
\end{description}

\tablerecord{\textbf{O14. A missed frame}}
\begin{description}
\item[{OrganoidTracker (where)}] \texttt{links\_\allowbreak{}postprocessor.\allowbreak{}bridge\_\allowbreak{}gaps}
\item[{what it does there}] A loose end at t and a loose start at t+2, each among the other's 6 nearest within 10 µm and neither with a nearer loose partner. If 2 < disappearance + appearance, a point goes in at t+1, at their midpoint. A point with a strong division score and one future counts as a loose end too.
\item[{ours: exact, nothing chosen}] S6's gap bridge: the end's prediction is carried across any number of frames; the pair is mutual nearest; and the inserted point is the exact midpoint, held as a sum and a count and written as a half voxel where odd. No distance or count is chosen.
\item[{stage}] S6
\item[{status}] theory
\end{description}

\tablerecord{\textbf{O15. A duplicate in one frame}}
\begin{description}
\item[{OrganoidTracker (where)}] \texttt{links\_\allowbreak{}postprocessor.\allowbreak{}bridge\_\allowbreak{}gaps2}
\item[{what it does there}] A loose end and a loose start in the same frame, within 7 µm and each the other's nearest. The end is a duplicate: its past is linked to the start, and it is removed.
\item[{ours: exact, nothing chosen}] Two points of one frame, one ending and one starting, mutual nearest: one body with two tops.
\item[{stage}] S6
\item[{status}] theory
\end{description}

\tablerecord{\textbf{O16. The division's frame}}
\begin{description}
\item[{OrganoidTracker (where)}] \texttt{links\_\allowbreak{}postprocessor.\allowbreak{}pinpoint\_\allowbreak{}divisions}, \texttt{remove\_\allowbreak{}spurs\_\allowbreak{}division}
\item[{what it does there}] A division moves one frame later when a daughter's division penalty is at least 1 below the mother's. A daughter that ends at once is deleted when the mother's division penalty is positive.
\item[{ours: exact, nothing chosen}] S7 puts the split at the frame of the strongest evidence along the chain. A daughter chain that does not stand above the null is no daughter.
\item[{stage}] S7
\item[{status}] theory
\end{description}

\tablerecord{\textbf{O17. Clean-up}}
\begin{description}
\item[{OrganoidTracker (where)}] \texttt{remove\_\allowbreak{}tracks\_\allowbreak{}too\_\allowbreak{}deep} (both ends past z = 25), \texttt{remove\_\allowbreak{}tracks\_\allowbreak{}too\_\allowbreak{}short} (6 points or fewer, not reaching the last frame), \texttt{remove\_\allowbreak{}single\_\allowbreak{}positions}
\item[{what it does there}] Removes the tracks and points so described.
\item[{ours: exact, nothing chosen}] Not taken as written: every length, depth and margin there is chosen. A chain is kept where it stands above its null (S4, S6).
\item[{stage}] S6
\item[{status}] not taken
\end{description}

\tablerecord{\textbf{O18. Marginals}}
\begin{description}
\item[{OrganoidTracker (where)}] \texttt{local\_\allowbreak{}marginalization\_\allowbreak{}functions}; \texttt{organoid\_\allowbreak{}tracker\_\allowbreak{}callibrate\_\allowbreak{}marginalization.\allowbreak{}py}
\item[{what it does there}] For each link, every flow-consistent state of the subgraph within 3 hops is weighted 10\textasciicircum{}−E, and P = Z\_\allowbreak{}link / Z. The hop count drops while the estimated count of states exceeds 2¹⁶. Options that leave the subgraph are folded into the node's end (\texttt{combine\_\allowbreak{}events}, 1 − Π(1 − p)). A division enters as a node's second unit. A temperature of 1.5, fitted by Platt scaling against curated links, scales the energies. Links with |penalty| ≥ 4 use only the target's in-links.
\item[{ours: exact, nothing chosen}] Built as the machine's marginal (M6) with the program's weights (row 7's \texttt{-\allowbreak{}-\allowbreak{}marginal}): exact integers, the masses cancel, no temperature. Over the new driver's points, the local set is the support component (low fruit 2), so nothing is outside and nothing is folded. Divisions are two-target arrangements (low fruit 3).
\item[{stage}] S6
\item[{status}] the marginal built (M6); over points, not built
\end{description}

\tablerecord{\textbf{O19. Dropping uncertain links}}
\begin{description}
\item[{OrganoidTracker (where)}] \texttt{organoid\_\allowbreak{}tracker\_\allowbreak{}marginalization.\allowbreak{}py}
\item[{what it does there}] A link with a marginal under 0.99 is dropped, and with it every out-link of a mother. Where the division score is high and no division was made, a second daughter is added at (x + 1, y + 1).
\item[{ours: exact, nothing chosen}] What is kept is each source's heaviest outcome in its marginal (an argmax, no threshold). No point is invented.
\item[{stage}] S6, S11
\item[{status}] theory
\end{description}

\tablerecord{\textbf{O20. Error flags}}
\begin{description}
\item[{OrganoidTracker (where)}] \texttt{cell\_\allowbreak{}error\_\allowbreak{}finder}, \texttt{errors.\allowbreak{}py}
\item[{what it does there}] Each point is checked for: a track end before the last frame; no past after the first; more than two daughters; a merge (two pasts); a low division or link score; a short cell cycle; a missed division (a high score, one future, and not high in the next two frames); moving too fast.
\item[{ours: exact, nothing chosen}] The graph's own invariants, checked before the output: in-degree ≤ 1; out-degree ≤ 2, and 2 only with a division verdict; every end a death, a leaving or the last frame; every start an entering, a daughter or the first frame. They are truthy records, counted per sample. The speed and cycle limits are chosen, so they are read from the set's distributions, not written in.
\item[{stage}] S11
\item[{status}] not built
\end{description}

\tablerecord{\textbf{O21. Local flow}}
\begin{description}
\item[{OrganoidTracker (where)}] \texttt{particle\_\allowbreak{}flow\_\allowbreak{}calculator}
\item[{what it does there}] The mean displacement of the neighbours within ±50 px in x and y and ±2 planes in z, each to its single next or previous point, skipping divisions and ends.
\item[{ours: exact, nothing chosen}] S5's local flow: the neighbours' departures from the drift, each weighted by its level. The neighbours are the point's neighbours in the Delaunay graph of the frame's points under ‖·‖\_\allowbreak{}w, with no radius, and divisions and ends are skipped as there. The first-pass links give it, and S10 sorts again.
\item[{stage}] S5, S10
\item[{status}] theory
\end{description}

\tablerecord{\textbf{O22. Image offsets}}
\begin{description}
\item[{OrganoidTracker (where)}] \texttt{core/\allowbreak{}images.\allowbreak{}Image\allowbreak{}Offsets}, set by hand in the offset editor; \texttt{is\_\allowbreak{}inside\_\allowbreak{}image}
\item[{what it does there}] A per-frame offset of the image. Positions are kept in its frame, and in-view tests read it.
\item[{ours: exact, nothing chosen}] S2's drift is the offset, read from the data per frame pair (E3).
\item[{stage}] S2, S9
\item[{status}] built in the engine
\end{description}

\subsection{\texorpdfstring{Hawkins et al. 2025, into the driver}{Hawkins et al. 2025, into the driver}}

\textbf{Source.} Hawkins, Jones, Wesselman, Cesa, Moeller, Wingert, “Mathematical modeling reveals cell differentiation processes and progenitor kinetics necessary for proper nephrogenesis”, \emph{Front. Cell Dev. Biol.} 13:1695380 (11 December 2025). Doug's upload, read 25 September.

\textbf{The model.} ODE models of the zebrafish pronephros from 10 to 96 hpf. Two progenitor pools (anterior G, posterior J) feed three tubule populations (anterior A, posterior P, multiciliated M). Every rate (proliferation, fate, death) is a Gaussian in time, K(t) = K\_\allowbreak{}max · 2\textasciicircum{}(−((t − K\_\allowbreak{}time)/K\_\allowbreak{}width)²). The models are fitted by bounded BFGS, simulated annealing and Monte Carlo over 100 × 100 grids of starting counts, and compared by AIC and likelihood-ratio tests.

\textbf{The measurements.}

\begin{itemize}
\item PH3 (G2/M): about 7.5\% of PST cells and 3\% of DE cells at 20 hpf.
\item Caspase-3 (apoptosis): near 0 in both.
\item Zebrafish cycles from the literature: 6 to 10 h, with G2/M about 1 h, in retina and hindbrain at 28 hpf; 2 to 24 h for progenitors, by tissue and stage.
\end{itemize}

\textbf{The finding.} The best model needs three fates from a progenitor:

\begin{itemize}
\item symmetric division (two progenitors);
\item asymmetric division (one progenitor and one mature cell);
\item a fate change without division (the cell matures without dividing).
\end{itemize}

What it gives the tracker is physics, with none of its numbers written in:

{\footnotesize
\setlength{\tabcolsep}{6pt}
\begin{longtable}{>{\RaggedRight\arraybackslash}p{\dimexpr0.1246\linewidth-2\tabcolsep\relax}>{\RaggedRight\arraybackslash}p{\dimexpr0.2789\linewidth-2\tabcolsep\relax}>{\RaggedRight\arraybackslash}p{\dimexpr0.4824\linewidth-2\tabcolsep\relax}>{\RaggedRight\arraybackslash}p{\dimexpr0.1141\linewidth-2\tabcolsep\relax}}
\toprule
\# & the physics & what it changes in the driver & stage \\
\midrule
\endhead
\textbf{H1} & Asymmetric division: the two daughters are different cells. & A division does not need its daughters alike. Mass is conserved on their sum (row 9's D\_\allowbreak{}mass already reads β(m\_\allowbreak{}c + m\_\allowbreak{}s) = β(m\_\allowbreak{}p)), and they part along the parent's axis. Their shapes (S1b) and levels may differ, and from then on each daughter's chain carries its own harmonics. & S7 \\
\textbf{H2} & A fate change without division: a cell changes what it is without dividing. & A chain goes on through a change of shape or level with no split, because identity is coherence (row 7), not a constant look. S5 weighs a shape against the chain's last point, never its first, so a gradual change is not a break. A change of harmonics with no second point in the gate is a fate change, not a new body. & S5, S6 \\
\textbf{H3} & Death is near zero in the early pronephros (Caspase-3 about 0 in PST and DE at 20 hpf), and a few percent of cells are in G2/M at a time. & This agrees with S8: a chain that ends away from a face, with no bump, is a missed detection, not a death. Divisions are rare against continuations, as the keys show (87 of 199 hold any division, at most 5). The rates are read from each sample's own counts after the sort, never written in. & S8, S7 \\
\textbf{H4} & The population balances: its change is births minus deaths plus arrivals. & On any output with in-degree ≤ 1 and out-degree ≤ 2, N(t+1) − N(t) = D(t) − E(t) + S(t+1) exactly. D counts the two-daughter points of t, E the points of t with no future, and S the points of t+1 with no past. The physics says every end is a death, a leaving or the last frame, and every start is an entering or the first frame. The ends and starts left unexplained are the verdict, counted per frame and per sample. & S11 \\
\textbf{H5} & Each process's rate varies in time and peaks once (their Gaussian K(t)). & The sample's division count per frame is read as a curve over the whole sample after the sort: a set-wide reading beside S4's, with no shape fitted to it. & S4, S7 \\
\textbf{H6} & Hypotheses are compared by likelihood ratio (AIC, LRT). & Ours is the null (S4). Each event (a division, two tops, a missed frame) is kept where its exact count stands above the same count in the null frames. No χ² and no parameter penalty. & S4, S6 \\
\bottomrule
\end{longtable}
}

\section{\texorpdfstring{The table}{The table}}

\tablerecord{\textbf{1. The bodies}}
\begin{description}
\item[{the physics we hold it to}] A body is what stands above the medium, whole; the field carves its boundary, and no gradient is tuned.
\item[{does today}] The program takes the machine's component tree (M3) and one cut: the level with the most components, proved against the forest's count. \texttt{grow\_\allowbreak{}leaves} makes the leaves. \texttt{group\_\allowbreak{}objects} joins leaves into objects by touch in one frame, transitively; on today's bodies nothing touches, so every object is one body (row 6). Since the split (23 September), the slide and the overlap are machine operators over probes and probe links (M3). The program turns key cells and divisions into probes and links and grades the census itself, in \texttt{cell\_\allowbreak{}tracking/\allowbreak{}src/\allowbreak{}slide/\allowbreak{}slide\_\allowbreak{}score.\allowbreak{}cu}. That file is rewritten to the probe interface and is \textbf{not yet compiled or hooked in}: it hooks into score\_\allowbreak{}sample's frame loop, and score\_\allowbreak{}sample has to move to cell\_\allowbreak{}tracking first.
\item[{wants}] One node per cell. 44b6 is estimated at about 260 cells a frame and 6bba at about 700; the tree gives far more pieces.
\item[{tried, and what it gave}] Every node: \textbf{measured}, 5.8 M nodes against 880,906 estimated, score 0.279. The 400 largest a frame: \textbf{measured}, 0.661 on the 25 44b6. \texttt{stands}: \textbf{measured}, 0.12 on the older dump (38,114 nodes kept). \texttt{merge\_\allowbreak{}split}: on by default, no A/B. \texttt{dish}, \texttt{cohere}, \texttt{accrue}, \texttt{agree}: no run in the ledger. The slide (\texttt{-\allowbreak{}-\allowbreak{}slide}, 23 September, before the split), \textbf{measured}. Every level where the key cells' sharing changes is found by bisection over the levels. Each probe's component count is proved against the tree's count at that level: 0 of 26,454 differ. On 6bba\_\allowbreak{}48816121, the cut holds 931 of 935 key cells, and 50 are alone on their leaf. At each frame's best single level, 903 are alone; at some level of their own chain, 925. The most-components level sits low in the tree: many small pieces, with the cells merged into a few giant blobs. Frame 0: the cut is at level 84,314 of 1,135,213; all 4 key cells are alone from 1,045,988 (204 components) down to 463,871. The overlap at the same residual (\texttt{-\allowbreak{}-\allowbreak{}overlap}, 23 September), \textbf{measured}. Frames t and t+1 are cut at the same exact residual r. At the drift, each component's distinct partners are counted from the exact (earlier root, later root) pairs; mutual is the AND, both ends with exactly one. The sampled r values are each frame's key-partition events and its cut. On 6bba\_\allowbreak{}48816121 (99 pairs, 3,615 cuts), the earlier frame's cut has 61,679 components, of which 15,979 overlap exactly one and 6,688 are mutual (11\%). The best mutual over each pair's cuts sums to 13,783. Mutual peaks high in the tree, at 150 to 200 components a frame with 70–80\% of them mutual. Frame 0→1 at level 1,045,988: 204 components and 144 mutual; at the cut: 329 and 29. On 44b6\_\allowbreak{}0113de3b, the cut has 39,507 components and 7,962 mutual (20\%). The event census (23 September), \textbf{measured}. It classes every component: 1→1 move, 1→2, 2→1, 1→0 and 0→1. Mutual counted from either side agrees on every cut. On 6bba\_\allowbreak{}48816121 at the cut, about 70\% of the components have no partner at all. Of the key's 5 divisions, 0 are parted at the most-components cut. At each pair's best sampled cut, 5 are present, 3 have their daughters apart, 2 are parted (the parent's component overlaps both daughters' components) and 1 is caught as 1→2. Daughters separate only where the parent is gone, and one global drift misses daughters that move apart. Key cells per leaf (\texttt{-\allowbreak{}-\allowbreak{}box-\allowbreak{}history}, 23 September), \textbf{measured}. On 6bba\_\allowbreak{}48816121, 881 of the 931 key nodes that land on a leaf share it with another key node, on 110 leaves. Up to 13 key cells sit on one leaf of 779,789 voxels, and every key division's parent sits on a leaf of 652,575 to 1,004,159 voxels whose box is the whole 64×256×256 view. On 44b6\_\allowbreak{}0113de3b, 0 of 52 share. So ℓ* cuts 6bba into a few giant components that hold many cells each. The 6bba edge scores (856/918 on this sample) count leaf-to-leaf links, and most of those leaves hold many cells, so they are not per-cell tracking.
\item[{status}] built; the node count is the score's lever; on 6bba the cut holds many cells a leaf
\item[{next}] Bodies, not pieces, as the nodes: row 2's fingerprint and row 9's entropy decide membership. Hook \texttt{slide\_\allowbreak{}score.\allowbreak{}cu} in once score\_\allowbreak{}sample moves. The machine's LCA identity (M3, theory) gives every level where two key cells part without bisection: one query per neighbouring pair of key cells.
\end{description}

\tablerecord{\textbf{2. Each body's intrinsic state: its spherical harmonics}}
\begin{description}
\item[{the physics we hold it to}] Every body has unique harmonics because of intrinsic physical differences (mass density, membrane elasticity, aspect ratio). Cells are oblate, so ℓ = 1 (dipole) and ℓ = 2 (quadrupole) dominate. Up to ℓ = 2 the harmonics carry exactly what the moments carry: ℓ = 0 the mass, ℓ = 1 the index sums, ℓ = 2 the six second moments.
\item[{does today}] The machine accumulates every body's exact moments (M4). \texttt{grow} carries them to \texttt{Tree\allowbreak{}Frame.\allowbreak{}sums} and \texttt{.\allowbreak{}moments}. \texttt{coherence} exports them to the viewer. The print (\texttt{fingerprint\_\allowbreak{}program}, 52 steps) writes each body's ℓ ≤ 2 fingerprint as a record: the mass band and six signed golden bands of K′/m², K′ = sₐs\_\allowbreak{}b(m·M − SₐS\_\allowbreak{}b) scaled by \texttt{voxel\_\allowbreak{}pm} over its common unit. \texttt{print\_\allowbreak{}pair} weighs two prints against each other. \textbf{The linker (E8 in the equation below) does not read the print yet.} The sums are also read by \texttt{parallax}, \texttt{arc} (\texttt{track\_\allowbreak{}bend}), \texttt{velocity}, \texttt{division} and the export.
\item[{wants}] The print in the linker's pool, to tell a cell from its neighbour. ℓ > 2 needs the body's surface: theory. Oblate spheroidal harmonics in place of spherical: theory.
\item[{tried, and what it gave}] The print: \textbf{measured}, 81,807 bodies of 44b6\_\allowbreak{}0113de3b in 11.7 ms, CRC-64 \texttt{590f3a2c386647e3} twice, 0 differ from the host. Linking by it: see row 7 (\texttt{-\allowbreak{}-\allowbreak{}print-\allowbreak{}match}).
\item[{status}] built, read only by \texttt{-\allowbreak{}-\allowbreak{}print-\allowbreak{}match}
\item[{next}] Put the print into E8's pool: waits on Doug (build plan, Waiting on Doug 3).
\end{description}

\tablerecord{\textbf{3. The medium: the fluid the bodies sit in}}
\begin{description}
\item[{the physics we hold it to}] The wide term is the medium's scattering, its turbidity, the fluid's mean entropic flux. It is noise to us and is subtracted whole, in exact two's complement, so the medium resolves to zero (phase cancellation).
\item[{does today}] The program hands the machine's residual (M2) one smooth order and one background order per axis, \{2, 34, 34\} and \{6, 96, 96\} (\texttt{SMOOTH\_\allowbreak{}ORDERS}, \texttt{BACKGROUND\_\allowbreak{}ORDERS}, track.h). \textbf{They are a debt} (23 September): they came from \texttt{n = 4(σ/\allowbreak{}v)²} rounded to even, with two chosen scales, σ = 6/5 µm (smooth) and 2 µm (background) (\texttt{for\_\allowbreak{}anchor\_\allowbreak{}sift/\allowbreak{}stripped/\allowbreak{}src/\allowbreak{}exact\_\allowbreak{}track.\allowbreak{}py}). The rounding and the chosen σ both break the rules.
\item[{wants}] Orders read from the data, per sample, with nothing rounded. Doug's rulings, 23 September: the per-axis coherence reading gives an exact integer period P in voxel steps, and \textbf{order = period} (n = P). Flatten takes a smooth and background pair per sample (not built). The medium's physics now has an exact form (M2's heat identity): the residual is 4\textasciicircum{}m times what m explicit heat steps of the medium remove from the smoothed field, so “the medium's mean entropic flux” is literally the subtracted term.
\item[{tried, and what it gave}] The key against the passes, and the unit sweeps: see M2 (proved exact; the sweeps 2× faster), with its numbers on the samples in \hyperref[chap:on-the-engine]{on\_\allowbreak{}the\_\allowbreak{}engine.md}.
\item[{status}] orders chosen and rounded (debt)
\item[{next}] Build the per-axis coherence reading (machine, M2's next), then orders per sample. Whether the tracker runs the unit sweeps by default (\texttt{-\allowbreak{}-\allowbreak{}unit-\allowbreak{}sweep}) is Doug's call.
\end{description}

\tablerecord{\textbf{4. Microfluidics: the medium's flow}}
\begin{description}
\item[{the physics we hold it to}] The field moves as a whole (drift) and each body moves on top of it; the whole moves first. Mitosis and lysis change the fluidics for a great distance around them: a division or a lysis sends a hydrodynamic disturbance through the surrounding tissue. Oblate bodies shed predictable dipole waves when they deform, snap or throw out a pseudopod.
\item[{does today}] Bulk drift only: the machine's drift (M5, \texttt{shift\_\allowbreak{}agreement}) gives the one lag that best carries this frame's set bits onto the next frame's, \texttt{lag\_\allowbreak{}to\_\allowbreak{}next} (the demon's eyes). Nothing reads a local flow, a disturbance or a wave.
\item[{wants}] The disturbance field around each event read from the residual, so an event far off is felt where it arrives. Overlapping disturbances untangled by their harmonic coefficients, each traced back to its origin, retrograde.
\item[{tried, and what it gave}] \texttt{motion\_\allowbreak{}check} (drift run twice, every count compared): built, no disagreement recorded. \texttt{keep\_\allowbreak{}view} (no drift removed): no run in the ledger.
\item[{status}] drift built; flow theory
\item[{next}] Build plan C3: each body's departure from the drift (v*\_\allowbreak{}b − d\_\allowbreak{}t), summed over its contacts and weighted by mass, then traced back toward the division and lysis events it came from.
\end{description}

\tablerecord{\textbf{5. Each body's motion}}
\begin{description}
\item[{the physics we hold it to}] A body's velocity is the time derivative of its dipole. Its next place is predicted from that before any overlap is counted, and the dipole axis points back to where the motion started: v\_\allowbreak{}origin = −∇(∂C₁ₘ/∂t).
\item[{does today}] The climb (M5, \texttt{climb\_\allowbreak{}machine}): each body climbs from the drift to its own lag, one step in 26 directions at a time, by how many of its voxels agree, until no step helps. \texttt{-\allowbreak{}-\allowbreak{}velocity} writes the dipole's change per body as a record beside the climb (see the verdicts). \textbf{\texttt{Engine\allowbreak{}Body.\allowbreak{}velocity} exists and nothing writes it; the linker does not read the velocity record.}
\item[{wants}] Velocity from the dipole's change across frames, used to predict each landing, and checked against the climb.
\item[{tried, and what it gave}] \texttt{sticky}, \texttt{mass} (land where most of the body lands), \texttt{spiral} (golden spiral offsets beyond the climbed lag), \texttt{arms} (the same climb at gaps 2 to n, so bodies too close one frame on have parted later): all built, no run in the ledger. \texttt{-\allowbreak{}-\allowbreak{}velocity} (the \texttt{velocity} record program): the dipole's change is written per body as an exact rational, (S′m − Sm′)/(m·m′) per axis, beside a truthy verdict that the climbed lag stands on a lattice point next to it. \textbf{Measured} on 44b6\_\allowbreak{}0113de3b, device equal to the host on every lane: of 35,755 bodies with a climbed forward, the climb and the dipole agree on all three axes for 8,720 (z 22,660, y 13,227, x 13,306); the last change predicts the next climb for 752 of 35,966. On 6bba\_\allowbreak{}48816121, 6bba\_\allowbreak{}09961292 and 6bba\_\allowbreak{}cdcfe533, \textbf{measured}: the check agrees on all axes for 2,144 of 21,366, 22,027 of 60,714 and 5,466 of 29,252; the prediction for 201 of 21,724, 2,666 of 58,840 and 497 of 28,743; 0 differ from the host.
\item[{status}] climb built; dipole velocity built and measured on four samples
\item[{next}] The leaves are pieces (818 a frame against about 260 cells), so a piece's dipole moves with the cut, not the cell: measure the same on bodies, and on the 25.
\end{description}

\tablerecord{\textbf{6. Protein interactions: contact and adhesion}}
\begin{description}
\item[{the physics we hold it to}] Cells touch and adhere through membrane proteins; bodies in contact move together and pull on each other. The membrane is a scale of the problem (\texttt{membrane\_\allowbreak{}pm}, 4 nm, the bilayer's thickness).
\item[{does today}] \textbf{No two bodies touch on today's engine} (corrected 23 September). The bodies are the separate components of the cut at ℓ*, and two components cannot share a face without being one. So \texttt{grow\_\allowbreak{}leaves} holds no touching pairs (\texttt{joined\_\allowbreak{}count} = 0), and \texttt{group\_\allowbreak{}objects}, \texttt{sticky}, \texttt{web} and \texttt{-\allowbreak{}-\allowbreak{}contact-\allowbreak{}side} have nothing to read. The joined leaf pairs came from the older basin engine. The volume faces each body touches (\texttt{touches}) are recorded. \texttt{voxel\_\allowbreak{}pm} is read by the print, \texttt{division} and \texttt{contact\_\allowbreak{}side}. \texttt{membrane\_\allowbreak{}pm} and \texttt{species} are read into the \texttt{.\allowbreak{}cfg} and written back out, \textbf{with no reader}.
\item[{wants}] Contact as a force-free constraint: bodies adhered in one frame stay adhered unless the data parts them. The membrane scale in place of any voxel count. Membership never by touch alone (row 9).
\item[{tried, and what it gave}] \texttt{sticky}, \texttt{web}: built, no run in the ledger. Grouping by touch: \textbf{measured} on the older basin engine to chain up to 2,256 pieces into one object at the widest; on today's engine every object is one body. \texttt{-\allowbreak{}-\allowbreak{}contact-\allowbreak{}side} on 44b6\_\allowbreak{}0113de3b and three 6bba: \textbf{measured}, 0 touching pairs on every sample, because there are none.
\item[{status}] contacts not held on today's bodies; adhesion theory
\item[{next}] What “touching” is for components of a cut is Doug's to define: for example, the component tree's next level down, where two bodies first join, or a gap of the membrane's width. The tree gives the first of these exactly: two bodies first join at the level of their LCA (M3). Until then the no-crossing verdict (built, \texttt{-\allowbreak{}-\allowbreak{}contact-\allowbreak{}side}) has no pairs. B3's deformation waits on C1: the membrane deforms, so a body and its match overlap only in part. What stays inside the match under every jitter is the body's internal constituents, and their shift is the velocity and vector.
\end{description}

\tablerecord{\textbf{7. Correspondence: which body it is next frame}}
\begin{description}
\item[{the physics we hold it to}] Identity is the limit of coherence: a body is who it is because its coherence carries on, not because of a tag. One-to-one where bodies are one-to-one.
\item[{does today}] The overlap triples (M5, \texttt{body\_\allowbreak{}overlap}) count shared positive voxels per pair at the drift. \texttt{pick} keeps the branch sharing the most, \texttt{unbound} takes the nearest centre by the swept vector (the jitter step it meets at over its magnitude), and \texttt{resolve} rejoins reconverging branches by union-find. The machine's one-to-one matching (M7, \texttt{heaviest\_\allowbreak{}matching}) is wired in by \texttt{-\allowbreak{}-\allowbreak{}print-\allowbreak{}match}: the \texttt{print\_\allowbreak{}pair} record program weighs every candidate (the overlap triples and the climbed forward) by how far its seven-band print lies from the body's, and the heaviest one-to-one set rewrites the forwards. \texttt{-\allowbreak{}-\allowbreak{}marginal} writes each body's exact context-aware probability P(b → b′ | G) = Z\_\allowbreak{}\{b→b′\}/Z over its local set, with the machine's marginal (M6) and the program's weights (see the verdicts). It reads the triples and the null, and the linker does not read it yet. \textbf{Measured} (internal count): 98.4\% of 6,358 edges on the 25 44b6, 97.0\% of 3,873 on five 6bba.
\item[{wants}] Link each body to the one whose fingerprint agrees (row 2), one-to-one by \texttt{heaviest\_\allowbreak{}matching} on fingerprint cost.
\item[{tried, and what it gave}] Under the metric: \textbf{measured}, 87.1\% hit on 44b6 (10.9\% land on a neighbour, 5.75 µm) and 78.5\% on 6bba (19.6\%, 6.7 µm). \texttt{-\allowbreak{}-\allowbreak{}print-\allowbreak{}match} on 44b6\_\allowbreak{}0113de3b: \textbf{measured}, 56,927 pairs swept at once in 217 to 245 us, 31,130 of 80,697 leaves matched one to one, 19,113 forwards changed, edges unchanged at 38/12 of 50, because base.cfg's \texttt{unbound} linker picks by overlap weight and centre cost and reads the forward only as one more candidate. \texttt{-\allowbreak{}-\allowbreak{}marginal} on 44b6\_\allowbreak{}0113de3b (with the floor's null): \textbf{measured}, 34,297 bodies asked in 50 to 58 ms (30,444 in their context, 3,853 alone, 8 too wide), 123,538 options, device equal to the host. Of 49 key edges asked, the most probable in context is the truth for 16 and “no link” for 27. The heaviest candidate alone is the truth for 38, and the climb's forward for 35. Where a target wins, it is the truth 16 times of 22. The null's weight dominates: it is the null climb's held count, which the climb maximizes, while T is counted at the drift. Putting the two on one footing is Doug's to set. With Doug's transform (\texttt{-\allowbreak{}-\allowbreak{}box -\allowbreak{}-\allowbreak{}marginal}, the links weighed by T*, C at each body's own climbed lag): \textbf{measured}, 51,214 bodies asked in 240 ms, 195,658 options, device equal to the host. Of 49 key edges, the most probable in context is the truth for 19 and “no link” for 24; the heaviest candidate alone gets 37 and the climb's forward 35. “No link” is weighed by the best of the floor's null draws (max\_\allowbreak{}k h⁰\_\allowbreak{}k). Whether “the null draw's fraction” means the best draw, each draw as its own outcome, or something else is Doug's. No run in the ledger for \texttt{share}, \texttt{cast}, \texttt{parallax}, \texttt{arc}, \texttt{settle}, \texttt{damp}, \texttt{vote}, \texttt{mutual}, \texttt{forest}, \texttt{merge\_\allowbreak{}target}, \texttt{forward\_\allowbreak{}only}, \texttt{focus} or \texttt{tower}.
\item[{status}] built; losses are neighbours
\item[{next}] Row 2's test. The marginal's local set G is one hop today (the sources competing for b's candidates, and their candidates). M6's factorization says the exact marginal over the whole frame equals the marginal over b's connected component of the support graph, and over nothing smaller: take G as that component, and the probability is the global one exactly (theory).
\end{description}

\tablerecord{\textbf{8. The null: is a correspondence real}}
\begin{description}
\item[{the physics we hold it to}] A null is the field-noise reading, not an empty return: the same body climbed toward a frame where no correspondence can exist. A body is real where it stands above that reading. The arm is fired where the observer's cloud is densest, not swept.
\item[{does today}] Null draws (\texttt{-\allowbreak{}-\allowbreak{}null}), held against each body's own climb in the edge and node rows. With a floor laid down, there is one draw per floor on the identity line (frame \texttt{g·11 + f mod 11}), ordered by the window-overlap cloud (M8), least overlap first. \texttt{-\allowbreak{}-\allowbreak{}marginal} weighs each body's “no link” by its best null draw.
\item[{wants}] Nodes kept only where they stand above their null (\texttt{above\_\allowbreak{}null} in \texttt{maint/\allowbreak{}score\_\allowbreak{}submission.\allowbreak{}py}).
\item[{tried, and what it gave}] Aimed against swept, same samples and same count: to be measured. \texttt{above\_\allowbreak{}null} on the current engine, \textbf{measured} 23 September on the two samples with a floor (\texttt{-\allowbreak{}-\allowbreak{}null 1 -\allowbreak{}-\allowbreak{}run floor}, \texttt{maint/\allowbreak{}score\_\allowbreak{}submission.\allowbreak{}py}'s scorer: nodes matched by position within 7 µm, the score's node-count factor). 44b6\_\allowbreak{}0113de3b: all 81,807 nodes 0.027942 (edge jaccard 0.036, 2 of 50 key edges), \texttt{stands} 277 nodes 0, \texttt{above\_\allowbreak{}null} 7,890 nodes 0.039606 (jaccard 0.037). 6bba\_\allowbreak{}48816121: all 62,650 nodes 0.002323, \texttt{stands} 123 nodes 0, \texttt{above\_\allowbreak{}null} 1,012 nodes 0.001149. The tracker's own tally, by leaf identity, calls 38 of 44b6\_\allowbreak{}0113de3b's 50 key edges correct, where the position-matched scorer finds 2. The two measures disagree, and the submission is graded by position. Edge by edge on 44b6\_\allowbreak{}0113de3b: 11 of the 100 key-node ends have a centroid within 7 µm. The rest are 7 to 15 µm from the nearest one, because the key cell's leaf is a multi-cell blob (up to 24,481 voxels, about 3,000 for one cell) or a piece (3 to 357 voxels). The loss is the cut (row 1), not the link.
\item[{status}] built
\item[{next}] Run both on the 25: it is a node-count lever with no chosen number.
\end{description}

\tablerecord{\textbf{9. Division: mitosis}}
\begin{description}
\item[{the physics we hold it to}] A split is a boundary discontinuity, judged over the whole sample and not the moment. One body's fingerprint becomes two whose masses add to it and whose dipoles part along the parent's quadrupole axis. The local entropy bumps and settles into two lows.
\item[{does today}] \texttt{assign\_\allowbreak{}bodies} (only when node rows are asked for) marks a body SPLIT when nothing lands on it and its back landing is a body that went elsewhere, and gives it that body's id as parent. \texttt{-\allowbreak{}-\allowbreak{}division} writes the mass and axis verdicts for every parent and two pieces whose backward climbs land on it (the \texttt{division} record, 100 steps). \texttt{assign\_\allowbreak{}bodies} does not read them yet, and there is no entropy check.
\item[{wants}] Two linked children confirmed by mass conservation, the parent's ℓ = 2 axis, and an entropy bump that settles.
\item[{tried, and what it gave}] Divisions matched: \textbf{measured}, 0 on the older dump. The division term earns 0 today. \texttt{-\allowbreak{}-\allowbreak{}division} on 44b6\_\allowbreak{}0113de3b: \textbf{measured}, 142,649 triples in 13.8 ms, mass conserved on 3,356, split along the long axis on 99,155, both on 2,452, 0 differ from the host. That sample's key holds no division. Of the 199 keys, 87 hold any, at most 5 (6bba\_\allowbreak{}48816121). On 6bba\_\allowbreak{}48816121, 6bba\_\allowbreak{}09961292 and 6bba\_\allowbreak{}cdcfe533: \textbf{measured}, 326,836, 85,146 and 99,724 triples, each in 1 to 3 ms, 0 differ from the host. The keys hold 13 divisions. Only 3 put the parent and both children on three separate leaves, and 2 of those are among the triples. Mass holds on 1 of the 2, the axis on 2, both on 1. Why the rest are lost, \textbf{measured}: 10 of the 13 have both children on one leaf. One frame after the split, the cut at ℓ* has not separated the daughters. On 6bba\_\allowbreak{}48816121 (\texttt{-\allowbreak{}-\allowbreak{}box-\allowbreak{}history}), the parent and daughters of all 5 of its divisions are on leaves of 0.65 to 1.0 M voxels whose box is the whole view (row 1). None has a node off every leaf, and none falls outside consecutive tracked frames. On the 3 with three separate leaves, 4 of the 6 children climb back to their parent. So the loss is the cut (row 1), not the climb.
\item[{status}] fate built; mass and axis verdicts built and measured
\item[{next}] The daughters inside one leaf are likely separate components higher in the same component tree. One cut level for the whole frame cannot hold both a just-divided pair and everything else. That is row 1's selection (Ultrack's one node per nested chain; build plan, Waiting on Doug 5). The census's 1→2 class (M3) is exactly the set of sources a division arrangement would add to the marginal (row 7). Then build plan C2 (the entropy bump).
\end{description}

\tablerecord{\textbf{10. Death: lysis}}
\begin{description}
\item[{the physics we hold it to}] An uncontained dissipation: the entropy bump rises into the floor and does not settle, and the body bleeds into the background.
\item[{does today}] \texttt{assign\_\allowbreak{}bodies} marks VANISHED (nothing onward, not at a face), ABSORBED (landed in a body another carried) and MERGED. Nothing tells a lysis from a missed detection.
\item[{wants}] A lysis told from a vanishing by its bump rising into the floor.
\item[{tried, and what it gave}] The box history (C2's measurement): \textbf{proved} on 44b6\_\allowbreak{}0113de3b (81,807 boxes × 9 windows × 16 bits, 0 of 11,780,208 differ from walking every box, 5.4 s) and 6bba\_\allowbreak{}48816121 (0 of 9,021,600, 6.3 s). Every key division's parent on 6bba\_\allowbreak{}48816121 sits on a leaf whose box is the whole view (row 1), so its history is the view's, and a bump cannot be read there.
\item[{status}] measured, no verdict
\item[{next}] Build plan C2. \texttt{-\allowbreak{}-\allowbreak{}box-\allowbreak{}history} holds each body's box (\texttt{climb\_\allowbreak{}machine\_\allowbreak{}extents}) and gathers its flips per window and bit from one running-sum key a window and bit, by 8 corners. The verdict (a bump that rises and settles, or rises and stays) waits on row 1's cut and on Doug: which bits, and how the windows (11 transitions) meet a one-frame event.
\end{description}

\tablerecord{\textbf{11. Entering and leaving the volume}}
\begin{description}
\item[{the physics we hold it to}] Bodies cross the volume's faces; a body that appears at a face entered, one that disappears at a face left.
\item[{does today}] ENTERED and LEFT from \texttt{touches} (the volume faces a body meets).
\item[{wants}] Nothing further now.
\item[{tried, and what it gave}] none
\item[{status}] built
\item[{next}] none
\end{description}

\tablerecord{\textbf{12. The noise floor: what is not a body}}
\begin{description}
\item[{the physics we hold it to}] The noise is deterministic, constant and unique per sample: shot, thermal and read, fixed pattern, quantization. It is read, never modelled. Entropy separates it: what decays to maximum entropy is floor −4, and what holds low and moves is a body.
\item[{does today}] The machine's entropy history (M8) with its window-overlap cloud; the \texttt{.\allowbreak{}cfg}'s \texttt{floor} lays it down first.
\item[{wants}] Every departure from the floor placed, so only those places are asked about (the demon's waveform collapsing onto U | U). The four noise keys imprinted and stamped over the set in one cycle.
\item[{tried, and what it gave}] Anchors: \textbf{measured}, bits 0 to 4 in about 46\% of frames everywhere, anchors in bits 6 to 11, none common to every sample. History: \textbf{measured}, bits 0 to 3 at 499 to 500 per mille in 19 of 25. A constant region in 6 samples: measured, cause unknown. Direction: measured, each sample's own.
\item[{status}] measured
\item[{next}] Build plan C4 (waits on Doug's design), then membership by entropic direction (rows 1, 9, 10).
\end{description}

\tablerecord{\textbf{13. A body's edge}}
\begin{description}
\item[{the physics we hold it to}] The edge is what stays with the body under every jitter of the view; what falls in and out is the floor at its boundary, set aside and not averaged.
\item[{does today}] The jitter sweep (\texttt{LINK\_\allowbreak{}SWEEP\_\allowbreak{}STEPS}, a box a voxel wide doubling out past the view) exists and is used only to score a link's disagreement. The jitter core (M5, \texttt{climb\_\allowbreak{}machine\_\allowbreak{}core}, \texttt{-\allowbreak{}-\allowbreak{}core}) is built and measured.
\item[{wants}] Every boundary voxel scrubbed under the sweep, so a body's edge and mass are its own and not one frame's cut.
\item[{tried, and what it gave}] The core per width on 44b6\_\allowbreak{}0113de3b, \textbf{measured}. Proof: the width-0 core mass equals C at the climbed lag on 35,755 of 35,755 forward climbers. Of 161,708 climbers (both sides, adjacent frames), 72,186 reach width 0, 33,827 width 1, 22,884 width 2, 10,230 width 3, 606 width 4 and 2 width 5. Core mass is 90,205,526 / 36,098,925 / 6,542,528 / 307,434 / 4,956 / 2. Which axis ends each core is not yet measured. Mutual pairs whose core displacement lies within half a voxel of the climbed lag on every axis: 10,003 of 16,176 at w=0, 4,409/10,410 at w=1, 2,578/8,727 at w=2, 456/3,413 at w=3, 15/117 at w=4. The deeper the core, the less its motion is the lag's. 49.6 s, with D3 sharing the GPU.
\item[{status}] core built and measured
\item[{next}] Build plan C1 (Doug, 22 September: every width, kept per width, between frames): a voxel of a body stays at width w when every offset in the box of half-width w around the climbed lag lands it inside the match. Mass, sums and moments re-summed per width, so the print is summed over the edge that stays.
\end{description}

\tablerecord{\textbf{14. Size: growth and division by ratio}}
\begin{description}
\item[{the physics we hold it to}] Sizes spread on a golden ladder, since growth and division scale by ratios.
\item[{does today}] The machine's golden ladder (M11, \texttt{band\_\allowbreak{}of}), read by \texttt{damp}, by \texttt{group\_\allowbreak{}objects}, by the residual survey, by the print and by \texttt{division} (mass conserved is β(m\_\allowbreak{}c + m\_\allowbreak{}s) = β(m\_\allowbreak{}p)). \texttt{-\allowbreak{}-\allowbreak{}mass-\allowbreak{}band} (\texttt{band\_\allowbreak{}or\_\allowbreak{}count}) compares the bands instead of the counts wherever two sizes are compared: \texttt{group\_\allowbreak{}objects}, \texttt{arm\_\allowbreak{}landing}, \texttt{focus\_\allowbreak{}links}, \texttt{assign\_\allowbreak{}bodies} and \texttt{track\_\allowbreak{}bend}. Exact sums stay counts, since a band there would round (Doug, 22 September).
\item[{wants}] A size read on the ladder instead of any count.
\item[{tried, and what it gave}] \texttt{damp}: no run in the ledger. \texttt{-\allowbreak{}-\allowbreak{}mass-\allowbreak{}band}, D3 on the first 25 (23 September): identical to counts on every one of the 24 that ran both sides, with 6,134 edges and 5,176 correct / 227 wrong / 2 none / 729 missed. It is identical by construction: on the scored path every reader is gated off. \texttt{track\_\allowbreak{}bend}, \texttt{arm\_\allowbreak{}landing} and the scorer's member choice compare within one object's members, and every object is one leaf. \texttt{group\_\allowbreak{}objects} needs joined pairs, and there are 0. \texttt{focus\_\allowbreak{}links} needs \texttt{-\allowbreak{}-\allowbreak{}focus}; \texttt{assign\_\allowbreak{}bodies} needs nodes output. The 25th, 44b6\_\allowbreak{}551a5dba, loaded in the counts run and then failed its .iapx CRC in the bands run and on every read since, with the file's mtime unchanged (for Doug).
\item[{status}] built, inert
\item[{next}] Live only once objects hold more than one leaf, or once bodies touch (row 6).
\end{description}

\tablerecord{\textbf{15. The whole state, once}}
\begin{description}
\item[{the physics we hold it to}] The field's state is a set of bodies at every time, and one exact file holds all of it.
\item[{does today}] The machine holds the field (M9: \texttt{.\allowbreak{}iapx}, \texttt{.\allowbreak{}oapx}, \texttt{flatten}, \texttt{.\allowbreak{}bapx}). No file holds the bodies' fingerprints.
\item[{wants}] A file kind for the bodies and their ℓ ≤ 2 fingerprints, written once, so linking reads bodies and not voxels.
\item[{tried, and what it gave}] See M9, and \texttt{flatten}'s numbers on 44b6\_\allowbreak{}0113de3b in \hyperref[chap:on-the-engine]{on\_\allowbreak{}the\_\allowbreak{}engine.md}.
\item[{status}] field held by the machine; bodies' prints not held
\item[{next}] With row 2.
\end{description}

\section{\texorpdfstring{The tracker's equation}{The tracker's equation}}

The tracker, written out as the one function it computes, stage by stage, from the code as it runs under \texttt{cell\_\allowbreak{}tracking/\allowbreak{}base.\allowbreak{}cfg} (\texttt{pick}, \texttt{unbound}, \texttt{merge\_\allowbreak{}split}, \texttt{resolve}, \texttt{climb: machine}). Every quantity is an exact integer unless a line says otherwise. Each stage names the code, and the machine operator it reads (M-rows and A-sections of engine\_\allowbreak{}table.md), so the algebra can be worked on the formula and checked against the code.

\subsection{\texorpdfstring{Notation}{Notation}}

\begin{itemize}
\item t is a frame, x = (z, y, x) a voxel of the view Ω = [0, D) × [0, H) × [0, W), and I\_\allowbreak{}t(x) the frame's 16-bit reading.
\item w = (w\_\allowbreak{}z, w\_\allowbreak{}y, w\_\allowbreak{}x) = (16, 1, 1) are the axis weights (\texttt{AXIS\_\allowbreak{}WEIGHTS}, track.h). ‖v‖²\_\allowbreak{}w = 16v\_\allowbreak{}z² + v\_\allowbreak{}y² + v\_\allowbreak{}x². They equal σ∘σ, with σ = voxel\_\allowbreak{}pm / gcd(voxel\_\allowbreak{}pm) = (1,625,000, 406,250, 406,250) / 406,250 = (4, 1, 1), the same σ \texttt{contact\_\allowbreak{}side} and the print already scale by (see the low fruit below).
\item [·] is 1 when its statement holds and 0 otherwise.
\item β(n) is n's rung on the golden ladder (A-section A9).
\end{itemize}

\subsection{\texorpdfstring{E1. The residual (row 3): the machine's A1}{E1. The residual (row 3): the machine's A1}}

R\_\allowbreak{}t = 2\textasciicircum{}g · (B\_\allowbreak{}s ∗ I\_\allowbreak{}t) − B\_\allowbreak{}b ∗ (B\_\allowbreak{}s ∗ I\_\allowbreak{}t), with the program's orders s and b (row 3's debt). The positive set is P\_\allowbreak{}t(x) = [R\_\allowbreak{}t(x) > 0].

\subsection{\texorpdfstring{E2. The bodies (row 1): the machine's A2 and A3, then \texttt{grow\_\allowbreak{}leaves}}{E2. The bodies (row 1): the machine's A2 and A3, then grow\_leaves}}

\begin{itemize}
\item The cut is ℓ*\_\allowbreak{}t = min \{ ℓ : C\_\allowbreak{}t(ℓ) = max C\_\allowbreak{}t \}, the lowest level with the most components (A2). The choice of one level for the whole frame is the program's.
\item The bodies B\_\allowbreak{}t are the components at ℓ*\_\allowbreak{}t, with mass m\_\allowbreak{}b, index sums S\_\allowbreak{}b, second moments M\_\allowbreak{}b, peak p\_\allowbreak{}b and touches (A3). The centroid is carried as the pair (S\_\allowbreak{}b, m\_\allowbreak{}b) and never divided.
\end{itemize}

\subsection{\texorpdfstring{E3. The drift (row 4): A4, \texttt{shift\_\allowbreak{}agreement}}{E3. The drift (row 4): A4, shift\_agreement}}

d\_\allowbreak{}t = argmax\_\allowbreak{}v Σ\_\allowbreak{}x P\_\allowbreak{}t(x) · P\_\allowbreak{}\{t+1\}(x + v)

\subsection{\texorpdfstring{E4. The climb (row 5): A4, \texttt{climb\_\allowbreak{}machine}}{E4. The climb (row 5): A4, climb\_machine}}

For each body b of frame t, A\_\allowbreak{}b(v) = Σ\_\allowbreak{}\{x∈b\} P\_\allowbreak{}\{t+1\}(x + v). The climb starts at v₀ = d\_\allowbreak{}t, steps to the neighbour u (26 of them) with the largest A\_\allowbreak{}b(u) while A\_\allowbreak{}b(u) > A\_\allowbreak{}b(v\_\allowbreak{}k), ties to the least ‖u‖²\_\allowbreak{}w, and stops at v*\_\allowbreak{}b.

\begin{itemize}
\item The forward is f(b) = L\_\allowbreak{}\{t+1\}(p\_\allowbreak{}b + v*\_\allowbreak{}b), the body of t+1 under the peak, or none.
\item The climb's held score is h(b) = A\_\allowbreak{}b(v*\_\allowbreak{}b).
\item The backward runs the same from t+1 to t, starting at −d\_\allowbreak{}t, and gives back(b′) and its lag v*\_\allowbreak{}\{b′\}.
\end{itemize}

\subsection{\texorpdfstring{E5. The null (row 8): the same climb toward a frame where no correspondence exists}{E5. The null (row 8): the same climb toward a frame where no correspondence exists}}

h⁰\_\allowbreak{}k(b) = A\_\allowbreak{}b\textasciicircum{}\{(t → t\_\allowbreak{}k)\}(v*),  one draw k per floor on the identity line when a floor is laid down

\subsection{\texorpdfstring{E6. The overlap triples (row 7): A4, \texttt{body\_\allowbreak{}overlap}}{E6. The overlap triples (row 7): A4, body\_overlap}}

T(b, b′) = |\{ x : L\_\allowbreak{}t(x) = b, L\_\allowbreak{}\{t+1\}(x + d\_\allowbreak{}t) = b′ \}|,  kept where T > 0

\subsection{\texorpdfstring{E7. The objects (row 6): \texttt{group\_\allowbreak{}objects} (\texttt{merge\_\allowbreak{}split})}{E7. The objects (row 6): group\_objects (merge\_split)}}

b \textasciitilde{} b′ ⇔ b, b′ touch ∧ ( O\_\allowbreak{}\{t−1\}(back(b)) = O\_\allowbreak{}\{t−1\}(back(b′)) ∨ f(b) = f(b′) ≠ none )

The objects O\_\allowbreak{}t are the transitive closure of \textasciitilde{}, one frame after another. \textbf{On today's bodies, E7 changes nothing}: two components of one cut cannot share a face, so \texttt{grow\_\allowbreak{}leaves} writes no touching pairs and every object is one body (“largest 1” on every run).

\subsection{\texorpdfstring{E8. The link (row 7): \texttt{link\_\allowbreak{}objects\_\allowbreak{}unbound}}{E8. The link (row 7): link\_objects\_unbound}}

For an object O of t and each object Q of t+1 that O's triples or forwards reach:

W(O, Q) = Σ\_\allowbreak{}\{b∈O\} Σ\_\allowbreak{}\{b′∈Q\} T(b, b′)

A forward that reaches Q enters the pool with weight 0.

κ(b, b′) = min over the four pairings (lag or lead) of  ( s ≪ 40 ) | ( Σ\_\allowbreak{}a w\_\allowbreak{}a (ℓ\_\allowbreak{}a + ℓ′\_\allowbreak{}a)²  mod 2⁴⁰ )

\begin{itemize}
\item ℓ is b's forward lag and ℓ′ is b′'s backward lag, so ℓ + ℓ′ is how far the round trip misses.
\item s = min \{ j ≤ 12 : 2\textasciicircum{}j ≥ max\_\allowbreak{}a |ℓ\_\allowbreak{}a + ℓ′\_\allowbreak{}a| \} is the jitter step it meets at.
\item K(O, Q) is the least κ over the member pairs that reach Q.

Links(O) = \{ Q : W(O, Q) = max\_\allowbreak{}Q′ W(O, Q′) ∧ K(O, Q) = min over those Q of K \}
\end{itemize}

The link is the heaviest overlap. The round-trip cost only breaks ties.

\subsection{\texorpdfstring{E9. Reconvergence (row 7): \texttt{resolve}}{E9. Reconvergence (row 7): resolve}}

When O has two links whose single-link chains meet again at one node, every pair of nodes along the two chains is joined (union-find). The nodes are the classes that result.

\subsection{\texorpdfstring{E10. Fate (rows 9 to 11): \texttt{assign\_\allowbreak{}bodies}}{E10. Fate (rows 9 to 11): assign\_bodies}}

\begin{itemize}
\item SPLIT: nothing lands on the body, and its back landing is a body that went elsewhere. That body becomes its parent.
\item VANISHED, ABSORBED, MERGED, ENTERED and LEFT come from the landings and \texttt{touches}.
\end{itemize}

\subsection{\texorpdfstring{The verdicts built beside the path (records on the machine's record engine, M10; truthy: zero false, nonzero true; AND is a product, OR a sum)}{The verdicts built beside the path (records on the machine's record engine, M10; truthy: zero false, nonzero true; AND is a product, OR a sum)}}

These are written for every body, pair or triple in one sweep each, device equal to the host. \textbf{None of them enters E8 yet.}

\begin{itemize}
\item \textbf{The print} (row 2, \texttt{fingerprint}), per body: β(m) and the six signed golden bands of K′ / m², where K′ = s\_\allowbreak{}a s\_\allowbreak{}b (m M − S\_\allowbreak{}a S\_\allowbreak{}b) is scaled by \texttt{voxel\_\allowbreak{}pm} over its common unit.
\item \textbf{The print cost} (row 7, \texttt{print\_\allowbreak{}pair}), per pair: the distance between two prints. The ceiling 1,275 is a body against itself.
\item \textbf{Velocity} (row 5, \texttt{velocity}), per body and its forward, and per axis:

\begin{itemize}
\item Δ\_\allowbreak{}a = (S′\_\allowbreak{}a m − S\_\allowbreak{}a m′) / (m m′), held as the pair (numerator, m m′)
\item R\_\allowbreak{}a = L\_\allowbreak{}a · m m′ − (S′\_\allowbreak{}a m − S\_\allowbreak{}a m′)
\item c\_\allowbreak{}a = COMPARE(m m′, |R\_\allowbreak{}a|), and the verdict V\_\allowbreak{}a = c\_\allowbreak{}a (c\_\allowbreak{}a + 1). It is nonzero exactly when |Δ\_\allowbreak{}a − L\_\allowbreak{}a| < 1, the climbed lag standing on a lattice point beside the dipole's change.
\item V = V\_\allowbreak{}z V\_\allowbreak{}y V\_\allowbreak{}x.
\end{itemize}
\item \textbf{Division} (row 9, \texttt{division}), per parent p and two pieces c, s whose backward climbs land on p:

\begin{itemize}
\item D\_\allowbreak{}mass = 1 − |COMPARE(β(m\_\allowbreak{}c + m\_\allowbreak{}s), β(m\_\allowbreak{}p))|
\item D\_\allowbreak{}axis = [3 Δᵀ K\_\allowbreak{}p Δ ≥ tr(K\_\allowbreak{}p) |Δ|²], where Δ = S\_\allowbreak{}c m\_\allowbreak{}s − S\_\allowbreak{}s m\_\allowbreak{}c
\item D = D\_\allowbreak{}mass · D\_\allowbreak{}axis
\end{itemize}
\item \textbf{Contact side} (row 6, \texttt{contact\_\allowbreak{}side}, two passes), per touching pair A, B of frame t whose forwards A′, B′ both exist:

\begin{itemize}
\item pass 1: N = σ ∘ (S\_\allowbreak{}A m\_\allowbreak{}B − S\_\allowbreak{}B m\_\allowbreak{}A) per pair, where σ\_\allowbreak{}a = voxel\_\allowbreak{}pm\_\allowbreak{}a / gcd(voxel\_\allowbreak{}pm);
\item pass 2 reads two pass-1 records: dot = N · N′, c = COMPARE(dot + 1, 1) = sign(dot);
\item kept = c (c + 1) and crossed = c (c − 1).
\end{itemize}

dot is (c\_\allowbreak{}A − c\_\allowbreak{}B)·(c\_\allowbreak{}A′ − c\_\allowbreak{}B′) in physical space times the positive m\_\allowbreak{}A m\_\allowbreak{}B m\_\allowbreak{}A′ m\_\allowbreak{}B′ / gcd², so its sign is exact. Two membranes cannot pass through one another, so a real pair keeps its side. dot = 0 includes A′ = B′, the pair landing on one body.
\item \textbf{The marginal} (rows 7 and 8), the machine's A5 with the program's weights: ω(b″, b′) = T(b″, b′) (or T*, below) and ω(b″, none) = max\_\allowbreak{}k h⁰\_\allowbreak{}k(b″). P(b → b′ | G) = Z\_\allowbreak{}\{b→b′\} / Z. Every source's weights share its own mass as denominator, so the masses cancel, and Z and every Z\_\allowbreak{}\{b→·\} are exact integers compared by cross-multiplying, with no temperature, no energy and no calibration. OrganoidTracker 2.0's P(A | G) is the same sum over exp(−E/T) of calibrated network energies.
\end{itemize}

\subsection{\texorpdfstring{The whole, in one line}{The whole, in one line}}

Tracks = Resolve( Links( W(T(B(R(I)), d)), K(v*(B, d), v*(B′, −d)) ), Objects(B, f, back) )

This reads from the inside out:

\begin{enumerate}
\item the residual R;
\item its bodies B at the cut ℓ*;
\item the drift d between frames;
\item the overlap T at the drift and the climbed lags v*;
\item the objects from touch and shared landings;
\item the link, by the heaviest W and then the least K;
\item the reconverging chains joined.
\end{enumerate}

The records (print, velocity, division, marginal) are functions of the same B, v* and T. How they enter W and K is Doug's to set (Waiting on Doug, item 3, in the build plan).

\subsection{\texorpdfstring{The tracker's stages as reads of the machine's one operator (Doug, 23 September: “we are repeating the same operation close to 10 times”)}{The tracker's stages as reads of the machine's one operator (Doug, 23 September: "we are repeating the same operation close to 10 times")}}

Almost every stage above is the machine's correlation C (A4) followed by a different reduction:

{\small
\setlength{\tabcolsep}{6pt}
\begin{longtable}{>{\RaggedRight\arraybackslash}p{\dimexpr0.3570\linewidth-2\tabcolsep\relax}>{\RaggedRight\arraybackslash}p{\dimexpr0.6430\linewidth-2\tabcolsep\relax}}
\toprule
stage & the read of C \\
\midrule
\endhead
E2 moments m, S, M & \texttt{C\_\allowbreak{}\{t,t\}[a, a](0)} weighted by 1, x, x xᵀ \\
E3 drift d & argmax\_\allowbreak{}v Σ\_\allowbreak{}\{a,b\} \texttt{C\_\allowbreak{}\{t,t+1\}[a, b](v)} \\
E4 climb A\_\allowbreak{}b(v) & Σ\_\allowbreak{}\{b′ ∪ ∅\} \texttt{C\_\allowbreak{}\{t,t+1\}[b, b′](v)}, climbed to a local maximum v*\_\allowbreak{}b from d \\
E4 landing by mass & argmax\_\allowbreak{}\{b′\} C\_\allowbreak{}\{t,t+1\}b, b′ \\
the backward climb & the same box transposed: \texttt{C\_\allowbreak{}\{t+1,t\}[b′, b](−v)} = \texttt{C\_\allowbreak{}\{t,t+1\}[b, b′](v)} \\
E5 null h⁰ & Σ\_\allowbreak{}\{b′ ∪ ∅\} C\_\allowbreak{}\{t,t\_\allowbreak{}k\}b, b′, in a null frame \\
E6 overlap T & \texttt{C\_\allowbreak{}\{t,t+1\}[b, b′](d)} \\
the transform T* & C\_\allowbreak{}\{t,t+1\}b, b′, so Σ\_\allowbreak{}\{b′\} T* + T*(b, ∅) = h(b) \\
E8 link weight W & Σ over each object's members of T \\
C1 jitter core, every width & S\_\allowbreak{}w[b, b′] = Σ\_\allowbreak{}\{‖u‖∞ ≤ w\} C\_\allowbreak{}\{t,t+1\}[b, b′](v*\_\allowbreak{}b + u): nested boxes around v*\_\allowbreak{}b \\
contact κ\_\allowbreak{}w (row 6) & Σ\_\allowbreak{}v (B\_\allowbreak{}w ⋆ B\_\allowbreak{}w)(v) · \texttt{C\_\allowbreak{}\{t,t\}[A, B](v)}: the box's autocorrelation, a tent, over the frame's own box \\
the slide and the census (row 1) & C at pairs of levels, a rectangle sum of the own-node table (A2) \\
velocity, division, print & the moments above, read in pairs and triples \\
\bottomrule
\end{longtable}
}

What it replaces today: the drift computed once (\texttt{shift\_\allowbreak{}agreement}, over every v, by an exact NTT); the climb again per body (27 shifts a step); the overlap again at the drift; the landing again (\texttt{land\_\allowbreak{}mass}); the null again per draw; the core and contact each another pass; the census again per sampled cut. The machine's box (\texttt{-\allowbreak{}-\allowbreak{}box}) proved the climb, the null and the overlap are reads of the one C (engine table, M5; its numbers on 44b6\_\allowbreak{}0113de3b are in \hyperref[chap:on-the-engine]{on\_\allowbreak{}the\_\allowbreak{}engine.md}). The window of shifts is the open design choice, and it is the program's to name in its request: the machine holds no scale.

\subsection{\texorpdfstring{Doug's transforms, 23 September}{Doug's transforms, 23 September}}

\begin{itemize}
\item \textbf{The null's exact transform.} The link's weight is the climb's own functional split by the target it lands on:

T*(b, b′) = |\{ x ∈ b : P\_\allowbreak{}\{t+1\}(x + v*\_\allowbreak{}b) ∧ L\_\allowbreak{}\{t+1\}(x + v*\_\allowbreak{}b) = b′ \}|,  so Σ\_\allowbreak{}\{b′\} T*(b, b′) + T*(b, ∅) = A\_\allowbreak{}b(v*\_\allowbreak{}b) = h(b)

where ∅ is a positive voxel in no body (R > 0 below the cut). In the marginal, ω(b, b′) = T*(b, b′) and ω(b, none) = max\_\allowbreak{}k h⁰\_\allowbreak{}k(b). Both are counts of the same body's voxels under the same operator, and the masses still cancel.
\item \textbf{Contact} (row 6). Each body has an assigned boundary. At width w, κ\_\allowbreak{}w(A, B) = |\{ x : d\_\allowbreak{}A(x) ≤ w ∧ d\_\allowbreak{}B(x) ≤ w \}|, where d\_\allowbreak{}A(x) = min\_\allowbreak{}\{y∈A\} ‖x − y‖∞. It is nonzero exactly where the two boundaries overlap at w: the truthy contact, its size the magnitude.
\item \textbf{The jitter core} (C1, rows 13 and 6). The match b′ = f(b) has D\_\allowbreak{}\{b′\}(y) = min\_\allowbreak{}\{z∉b′\} ‖y − z‖∞. A voxel x of b stays at width w when D\_\allowbreak{}\{b′\}(x + v*\_\allowbreak{}b) > w. Built and proved: the width-0 core is Cb, b′, exactly, on every forward climber of 44b6\_\allowbreak{}0113de3b. The widths end at 5 there; the data sets where they end, and the view bounds them at (shortest extent + 1)/2.
\item \textbf{The cut} (row 1). ℓ* slides through every level of the component tree, to the top, and each nested chain keeps the node that stays coherent between frames. One level for the whole frame is replaced by one node per chain.
\end{itemize}

\subsection{\texorpdfstring{Low fruit in the tracker's algebra (23 September)}{Low fruit in the tracker's algebra (23 September)}}

Each item is algebra on the equation above, exact, \textbf{theory until built}. What it replaces is named.

\begin{enumerate}
\item \textbf{One scale vector.} w = σ∘σ with σ = voxel\_\allowbreak{}pm / gcd(voxel\_\allowbreak{}pm): the link's norm (E4, E8), the print's scale and the contact side's σ are the same vector. \texttt{AXIS\_\allowbreak{}WEIGHTS} becomes derived, not written in, and any geometry (another microscope, the MRI's 0.3 mm against 3.5 mm) gets its weights exactly: the machine's \texttt{decimal\_\allowbreak{}double} (A0) turns a spacing's decimal string into an exact integer, and the gcd makes them coprime integers. Replaces the constant \texttt{\{16, 1, 1\}}.
\item \textbf{The marginal over the component.} The exact frame-wide marginal factors over the connected components of the support graph (A5). b's probability is exact when G is b's whole component, and today's one-hop G is exact only when the component is one hop wide. Replaces the one-hop set; the “too wide” bodies (8 on 44b6\_\allowbreak{}0113de3b) are the components past \texttt{MARGINAL\_\allowbreak{}ARRANGEMENTS\_\allowbreak{}MOST}, which the machine's permanent on the component (A5) measures instead of guessing.
\item \textbf{Divisions as arrangements.} A division is a source that takes two targets. The census's 1→2 class (A2) is exactly where that arrangement is needed, so the marginal gains a two-target outcome only on those sources, and the rest keeps today's injective sum.
\item \textbf{The slide by LCA.} Two key cells share a component at level r exactly when r ≤ level(LCA). Ordered by the tree's DFS order, k key cells need k − 1 LCA levels to give their partition at every level at once (A2). Replaces the bisection over levels.
\item \textbf{The overlap at every level from one count.} The census at every sampled cut is one own-node pair table summed over DFS ranges (A2). Replaces a recount per cut.
\item \textbf{Contact is an LCA level.} Row 6's “where two bodies first join” is level(LCA(A, B)), exact, with no gap to choose.
\end{enumerate}

\section{\texorpdfstring{Where the tracker is not yet exact, or holds a chosen number}{Where the tracker is not yet exact, or holds a chosen number}}

{\small
\setlength{\tabcolsep}{6pt}
\begin{longtable}{>{\RaggedRight\arraybackslash}p{\dimexpr0.2578\linewidth-2\tabcolsep\relax}>{\RaggedRight\arraybackslash}p{\dimexpr0.4199\linewidth-2\tabcolsep\relax}>{\RaggedRight\arraybackslash}p{\dimexpr0.3223\linewidth-2\tabcolsep\relax}}
\toprule
where & what & in code \\
\midrule
\endhead
the residual's orders & \{2, 34, 34\} and \{6, 96, 96\} come from two chosen scales (6/5 µm, 2 µm) and n = 4(σ/v)² \textbf{rounded to even}. Doug, 23 September: no rounding, anywhere; order = period from the per-axis coherence reading, per sample & \texttt{SMOOTH\_\allowbreak{}ORDERS}, \texttt{BACKGROUND\_\allowbreak{}ORDERS}, track.h; \texttt{exact\_\allowbreak{}track.\allowbreak{}py} \\
w = (16, 1, 1) & equals σ∘σ with σ = voxel\_\allowbreak{}pm / gcd(voxel\_\allowbreak{}pm) = (4, 1, 1) exactly, but it is written in as a constant, not derived from \texttt{voxel\_\allowbreak{}pm} & \texttt{AXIS\_\allowbreak{}WEIGHTS}, track.h \\
κ's lead & (mine ± behind/2) / behind rounds by truncating division when frames are more than one apart; it is exact on consecutive frames (behind = 1), which is every frame today & \texttt{leaf\_\allowbreak{}disagreement}, relate\_\allowbreak{}frames.cu \\
κ's magnitude & kept mod 2⁴⁰, so it wraps silently above 2⁴⁰ & \texttt{LINK\_\allowbreak{}MAGNITUDE\_\allowbreak{}BITS} 40 \\
s & the jitter steps stop at 12, a chosen number & \texttt{LINK\_\allowbreak{}SWEEP\_\allowbreak{}STEPS} 12 \\
the marginal's set & one frame step and one hop; the exact marginal needs the whole support component (low fruit 2). Every target is taken once, so a division is not an arrangement yet (low fruit 3). The bound of 2¹⁶ arrangements is OrganoidTracker's, not measured here & \texttt{MARGINAL\_\allowbreak{}ARRANGEMENTS\_\allowbreak{}MOST} \\
the null in the marginal & ω(b, none) is the best null draw's held count, a count of the body's voxels like T, but T is counted at the drift and h⁰ at the climbed null lag & \texttt{score\_\allowbreak{}best\_\allowbreak{}null}, score\_\allowbreak{}sample.cu \\
\texttt{cohere} (off in base.cfg) & a mean motion is divided by truncating division & group\_\allowbreak{}objects.cu \\
the cut & one level (the most components) for the whole frame & \texttt{max\_\allowbreak{}tree\_\allowbreak{}objects} as the program calls it \\
\texttt{RUN\_\allowbreak{}PART\_\allowbreak{}ROOM} & 16, not measured & track\_\allowbreak{}driver.cu \\
\bottomrule
\end{longtable}
}

\section{\texorpdfstring{What the program asks of the engine, and what it offers}{What the program asks of the engine, and what it offers}}

The interface through which the two tables optimize against each other.

{\small
\setlength{\tabcolsep}{6pt}
\begin{longtable}{>{\RaggedRight\arraybackslash}p{\dimexpr0.4115\linewidth-2\tabcolsep\relax}>{\RaggedRight\arraybackslash}p{\dimexpr0.1710\linewidth-2\tabcolsep\relax}>{\RaggedRight\arraybackslash}p{\dimexpr0.4175\linewidth-2\tabcolsep\relax}}
\toprule
the program asks & the engine's row & state \\
\midrule
\endhead
probes and probe links, not key cells: the slide's partitions and the overlap's census & M3 & built (engine side); \texttt{slide\_\allowbreak{}score.\allowbreak{}cu} not hooked \\
the component tree's LCA levels and DFS order, to replace bisection (low fruit 4, 6) & M3 & theory \\
one own-node pair table at the drift, summed over DFS ranges (low fruit 5) & M3, M5 & theory \\
a per-axis coherence reading that returns an exact integer period (order = period) & M2 & theory \\
a smooth and background pair per sample in flatten & M9 & ruled, not built \\
the permanent over a whole support component, and its size & M6 & theory \\
exact spacings from decimal strings, for σ & M1 & built (\texttt{decimal\_\allowbreak{}double}) \\
a file for the bodies and their prints & M9 & theory \\
the residual over readings of any width (S0's contrast is the set's cumulative count, up to the set's reading count, far past 2¹⁶). The engine's integers are exact, bit-packed and of any width (Doug, 25 September). The request names the input's bits. The output's width is derived as input bits + Σ smooth orders + Σ background orders + 1 and returned, and peaks read that width & M2 & built 25 September. \texttt{unit\_\allowbreak{}sweep} takes limb planes of a declared width and refuses a value past it. \texttt{engine\_\allowbreak{}residual\_\allowbreak{}planes} returns the residual with its limbs, and \texttt{peaks\_\allowbreak{}find} takes the limbs. \textbf{Proved}: the planes against the 16-bit path, 336 checks, 0 failed (engine table, M2), and the peaks at four widths, 131 checks, 0 failed. The scan runs it at 34 bits in and 10 limbs out \\
\bottomrule
\end{longtable}
}

What the program offers the engine: nothing in the machine may know it. The measurements above are the engine's test load (44b6, 6bba), and the proofs the program runs (device equal to host, 0 lanes differ) are the machine's regression.

\section{\texorpdfstring{What to work toward, in order}{What to work toward, in order}}

The order is the build plan's (\hyperref[chap:build-plan]{build\_\allowbreak{}plan.md}), as of 23 September:

\begin{enumerate}
\item \textbf{Built as records beside the path:} the print (row 2); the print cost with one-to-one matching (row 7, \texttt{-\allowbreak{}-\allowbreak{}print-\allowbreak{}match}); velocity (row 5, \texttt{-\allowbreak{}-\allowbreak{}velocity}); division mass and axis (row 9, \texttt{-\allowbreak{}-\allowbreak{}division}); the exact marginal (rows 7 and 8, \texttt{-\allowbreak{}-\allowbreak{}marginal}); contact side (row 6, \texttt{-\allowbreak{}-\allowbreak{}contact-\allowbreak{}side}); sizes on the ladder (row 14, \texttt{-\allowbreak{}-\allowbreak{}mass-\allowbreak{}band}). None of them enters E8, the link, yet. How they enter W and K is Doug's to set.
\item \textbf{Move the cell-specific modules} to cell\_\allowbreak{}tracking: \texttt{score\_\allowbreak{}sample} and \texttt{coherence} moved 23 September, to \texttt{cell\_\allowbreak{}tracking/\allowbreak{}src/\allowbreak{}}, and \texttt{vis\_\allowbreak{}png} to \texttt{view/\allowbreak{}}. Then hook \texttt{slide\_\allowbreak{}score.\allowbreak{}cu} into the frame loop (row 1).
\item \textbf{Rows 6 and 13 (B3 with C1).} The jitter core per width.
\item \textbf{Rows 9 and 10 (C2).} The entropy bump that settles (division) or stays (lysis).
\item \textbf{Row 4 (C3).} Local flow from each body's departure from the drift.
\item \textbf{Row 3.} Orders per sample from the period reading, replacing the rounded debt; \textbf{row 12 (C4, waits on Doug).}
\item \textbf{D3.} Every want A/B on the 25, \texttt{above\_\allowbreak{}null} (row 8) included, written into the ledger and this table.
\end{enumerate}
